\chapter{How grammar feels} \label{ch:How grammar feels}

%\begin{quote}
% On the modest explanation, linguistic intuitions are, like any other intuitive judgments a person may make, immediate and unreflective reactions formed in a central processing system.
%\end{quote} \citep[4 (introduction)]{Schindler2020a}

%\begin{quote}
% ``I may long for certainty, but I have to live with doing the best I can. I may concoct a myth to explain that my certainty, unlike yours, taps into universal moral truths. Reality will soon dissolve that myth. Voltaire (1694–1778), a French philosopher of the Enlightenment, concisely summed up the state of affairs: \textit{Uncertainty is an uncomfortable position, but certainty is an absurd one.''\textsuperscript{10} Voltaire is right, of course, but nevertheless, we must make some choice or other, and not acting can be the worst of the options. Knowing that I am doing the best I can seems like cold comfort, yet my inner Socratic voice recommends that I make do with cold comfort.''\cite[Introduction]{churchland2019}
%\end{quote}

When you encounter an ungrammatical utterance, what does it feel like? Is it akin to stepping off the sidewalk unexpectedly~-- a sudden jolt to your cognitive equilibrium? Or sometimes more than just a simple misstep~-- a momentary disorientation, where the flow of language and thought breaks?

Is it sand in your shoe, a branch that scratches you as you walk by, an ice-cream headache, the sting of a mosquito, the glare of headlights? Does it feel heavy or light, salty or sour? Does it have a colour?

%Does ungrammaticality lean towards the cool side, like a draft that sneaks in through an unseen gap. Or is it a mouthful of scalding coffee? Is it a path that leads to a dead end or a door that's slightly ajar but won't fully open?

Emotionally, is this encounter more likely to evoke irritation and discomfort, fear and anger, or dismissal and superiority?

Or can it be all of those things at different times?

\bigskip

The idea that we have only five senses~-- sight, smell, hearing, taste, and touch~-- is nonsense. The traditional list of five ignores pain, pressure, itch, temperature, and chemical irritants like acids, the capsicum in hot peppers, and allyl isothiocyanate in radishes, wasabi and mustard. And that's just the exteroceptive senses~-- those that detect stimuli from outside the body.

We also have a rich array of interoceptive senses monitoring the internal states of our bodies, hunger for example, or sleepiness. When someone is dizzy, we literally say they've lost their sense of balance. Our sense of self relies on proprioception, the feeling we have of our bodies, their coherence, their locations, and their movements in space. We have senses of thirst, nausea, and fullness. We even have a sense of the levels of carbon dioxide and oxygen in our blood. We're sensitive creatures.

Sensations arise from the interaction between internal or external stimuli, our sense organs, and our nervous system. When our elbow is bent at 90$^{\circ}$, intrafusal fibers in our muscles register the relative amount of stretch in our biceps and triceps and Ruffini endings in the joint capsule and ligaments detect changes in joint angle. Nerves pick up the signals from these mechanoreceptors, and the sensations are sent as proprioceptive signals to our brain, informing it of the arm's position and movement. But positional sensations like this are typically evaluatively neutral. 90$^{\circ}$ should be neither more nor less pleasant than 120$^{\circ}$.

When we do evaluate sensations as good or bad, we've moved on to talking about \textsc{feelings}, the positive or negative experiences that come from sensing the world and ourselves. We sense hunger, pain, or high levels of CO\textsubscript{2} in the blood and we feel that they are unpleasant, problems to be righted. The sensation of warm touch on the arm, the sweetness of a peach, and the temporary relief that comes from scratching a mosquito bite, in contrast, foster positive feelings of relief, warmth, and connection, directing us to seek out such experiences.

\begin{quote}
    No sooner had the warm liquid mixed with the crumbs touched my palate than a shudder ran through me and I stopped, intent upon the extraordinary thing that was happening to me. An exquisite pleasure had invaded my senses, something isolated, detached, with no suggestion of its origin. And at once the vicissitudes of life had become indifferent to me, its disasters innocuous, its brevity illusory this new sensation having had on me the effect which love has of filling me with a precious essence; or rather this essence was not in me it was me. ... Whence did it come? What did it mean? How could I seize and apprehend it? ... And suddenly the memory revealed itself. \par\raggedleft~-- Marcel Proust, \textit{In search of lost time}
\end{quote}

Proust's vivid description highlights how a simple sensory encounter~-- the taste of a madeleine dipped in tea~-- can be intensely pleasurable. But the experiences transcends his sensory delight. He percieves a feeling whose source he can't identify at first, one triggered by but distinct from his experience of the smells, flavours, and textures of the moistened cake. Proust recognizes an upwelling associated with this particular set of sensations even before the memory arrives. It announces itself the way a yawn does, with a gradual build-up of tension before the mouth opens wide and the breath begins to move. A pricking of the thumbs. Summer's ripening breath.

This kind of feeling is called a metacognitive feeling. Unlike the taste of a madeleine or the warmth of a cup of tea, metacognitive feelings are generated by our own thoughts and mental operations. Such impressions include the positive feeling of knowing you might experience when reviewing your responses on a test, the pleasant recognition of a friend's familiar gate, or the empowering conviction of the rightness your opinion in a debate. It's not the thoughts that we have but the feeling that we have thoughts, whether we can access them or not. When Faulkner writes, ``memory believes before knowing remembers,'' he's evoking metacognitive feelings.

Metacognitive feelings are not all positive. Tip-of-the-tong feelings bedevil us. The thought that you've left something undone can be distressing. And then there's the feeling that this book is all about: ungrammaticality. This too I would place among our metacognitive feelings.


%Nick Chater, in \textit{The Mind is Flat}, argues that mental processes are fundamentally superficial, challenging the deep and complex models often posited in psychology and cognitive science. From this perspective, metacognitive feelings, including the sensation of (un)grammaticality, are immediate and surface-level responses, not rooted in any deep, structured understanding of grammar or cognitive processing. 

%Chater would likely suggest that the feeling of (un)grammaticality arises not from an intricate, rule-based analysis of language structures internalized through extensive learning, but from heuristic responses generated by the brain's superficial processing systems. These systems are quick to judge based on patterns and examples encountered in everyday language use, rather than deep grammatical rules. For instance, a sentence might feel ungrammatical simply because it's unfamiliar or deviates from frequently observed patterns, not necessarily because it violates a specific grammatical rule.

%This approach emphasizes how our perceptions of grammaticality are shaped more by pragmatic cues and less by any robust internal grammar engine. Thus, the feeling of (un)grammaticality is a product of our minds navigating through language based on surface impressions and immediate judgments, aligning well with Chater's view that the mind constructs its seemingly complex operations from moment to moment, without any underlying depth.

%\bigskip

%But we \emph{do} have a sense of grammar, or at least something it feels like to encounter grammaticality and ungrammaticality~-- not one thing, just like pain isn't one thing, but a set of related feelings that we experience when we encounter grammatical and ungrammatical sentences and which we associate with grammaticality and ungrammaticality.

When we identify a sentence as ungrammatical, our response is not primarily rooted in a physical sensation but in a cognitive evaluation. This process involves evaluating incoming linguistic constructions to ensure that they're familiar, conventional, and congruent, that they evoke the appropriate meaning for the situation, and that, if not, they deviate in understandable ways to good effect. It's difficult to identify necessary and sufficient criteria for feelings of ungrammaticality, but cognitive dissonance is often at their heart.

%This cognitive evaluation of grammaticality is reflective of our linguistic experience, a complex product of language exposure, learning, and use. It is through this lens that we can understand our reactions to ungrammatical constructions not as visceral responses but as cognitive acknowledgements of variance from the norm. The experience of encountering ungrammatical language, therefore, speaks to our capacity to detect, process, and interpret deviations within the framework of our linguistic knowledge.

%In this context, grammaticality is less about provoking a sensory reaction and more about the brain's ability to process language in accordance with learned patterns, what Dan Dennett called real patterns. %How do I connect this to chapter 10? Shall I just deal with it here?
Our reactions to grammatical incongruences may vary widely, from irritation to curiosity, reflecting the diversity in our linguistic backgrounds and tolerances for variation. Ultimately, the sense of grammar underscores the cognitive foundations of language processing and our adaptive strategies for navigating linguistic complexity.

%Is your sense of ungrammaticality the same feeling of wrongness you get when you see a stone levitating? I think it evokes more curiosity too, though. What about the sense of wrongness you feel when seeing unfairness?

%How do you react to your own grammatical slips?

%Perhaps it would be a step to far, though, to say we have a sense of grammar. I'll settle for saying that there is something it feels like to encounter an ungrammaticality~-- not one thing, just like pain isn't one thing, but a set of related feelings that we experience when we encounter ungrammatical sentences and which we associate with ungrammaticality.

\section{Whose gorilla?}

The article ``Whose gorilla?'', which appeared in \textit{Linguistic Inquiry} in \citeyear{Hankamer1973}, was a ``squib''~-- a short publication that presented an interesting piece of data without proposing an explanation. Not even running a full page, its punch line was (\ref{ex:gorilla}):

\ea \label{ex:gorilla}
 \textit{~My gorilla is over there drinking punch.} \\
 *\textit{The guy whose you saw banging at the window is over there watering}\\~\textit{the rubber tree.}
\z
They had fun in those days.

The authors of ``Whose gorilla?'' were Jorge Hankamer and Paul Postal. ``Postal was a venerable force, \dots~ second in reputation only to Chomsky in the prized domain of syntax'' \citep[127--128]{Harris2021}. He was also and outsider, spending most of his academic life at IBM instead of at a university, and something of a polemicist. His most famous paper was entitled ``The best theory'', and it supported a position called generative semantics that Chomsky would later deride as ``the worst theory, in fact as the worst possible theory \dots~the worst imaginable theory'' \citep[129]{Harris2021}. This despite the irony that ``the Chomsky of the 21st century reads like a reincarnation of the Postal of the 20th'' \citep[418]{Newmeyer2003}. 

As a graduate student, Postal was so sure of himself that he directed his own thesis. When he submitted it, his advisor wrote him a note saying ``he was too busy with his work to read it for two years'' \citep{PostalInterview2020}. This mixture of brilliance and confidence made him one of the key combatants in the so-called Linguistics Wars of the late `60s and early `70s, a truly rancorous battle between the generative semanticists on one side and Chomsky and his supporters on the other.

In 1973, Hankamer had just completed his PhD at Yale and taken up a position at Harvard. His dissertation, on deletion and ellipsis in Turkish, had already established him as a rising star in the field. Like Postal, he was on the outs with Chomsky despite being far less bombastic and pugilistic.

In their article ``Whose gorilla?'', Hankamer and Postal drew attention to an apparent gap in the pronoun system, specifically claiming that the word \textit{whose} is defective.

In linguistic terminology, a \textsc{defective} word is one that lacks one or more of the expected forms. This is why Hankamer and Postal have starred that sentence, to draw attention to the \textit{whose} that should~-- but seemingly couldn't~-- be there. 

What does it mean to lack a form? \textit{Beware}, for example is missing every form but one; it's highly defective. It has a bare form \textit{beware of the dog}, but no present-tense *\textit{she bewares the dog}, no past tense *\textit{She bewore the dog}, and no participles *\textit{I've never beworn the dog}.

This missing form~-- the purported hole in the \textit{whose} paradigm~-- should appear in the bottom row of Table \ref{tab:pronouns}, right where the question mark is.
%, though many linguists prefer the term \textsc{genitive}, which suggests a wider set of relationships than simple possession. When you speak of \uline{your} ancestors, \uline{your} birth, or \uline{your} situation~-- \textit{your} being a possessive/genitive pronoun~-- clearly you do not mean that you possess these things. A simple way to remember this is that the \uline{gen}itive suggests a \uline{gen}eral connection between things. Terminological issues aside, \textit{whose gorilla} evokes some kind of relationship with the gorilla.

\begin{table}[h]
\centering
\begin{tabular}{lccccc}
\hline
\textbf{Type} & \textbf{Nominative} & \textbf{Accusative} & \textbf{Dep Gen} & \textbf{Ind Gen} & \textbf{Reflexive} \\ 
\hline
\textbf{Personal} & \textit{I} & \textit{me} & \textit{my} & \textit{mine} & \textit{myself} \\ 
 & \textit{you} & \textit{you} & \textit{your} & \textit{yours} & \textit{yourself} \\ 
 & \textit{he} & \textit{him} & \textit{his} & \textit{his} & \textit{himself} \\ 
 & \textit{she} & \textit{her} & \textit{her} & \textit{hers} & \textit{herself} \\ 
 & \textit{it} & \textit{it} & \textit{its} & - & \textit{itself} \\ 
 & \textit{we} & \textit{us} & \textit{our} & \textit{ours} & \textit{ourselves} \\ 
 & \textit{they} & \textit{them} & \textit{their} & \textit{theirs} & \textit{themselves} \\ 
 & \textit{one} & \textit{one} & \textit{one's} & \textit{one's} & \textit{oneself} \\
\textbf{Interrog} & \textit{who} & \textit{whom} & \textit{whose} & \textit{whose} & - \\ 
\textbf{Relative} & \textit{who} & \textit{whom} & \textit{whose} & ? & - \\
\hline
\end{tabular}
\caption{Personal, Interrogative, and Relative Pronouns}
\label{tab:pronouns}
\end{table}

\textit{Whose} is often known as a \textsc{genitive pronoun}, AKA ``possessive pronoun''. The genitive pronouns come in two kinds: those that are \textsc{dependent} on a noun: \textit{\uline{my} life}, \textit{\uline{your} name}, \textit{\uline{her} eyes}, and those that are \textsc{independent}: \textit{mine}, \textit{yours}, \textit{hers}, etc. Perhaps confusingly, dependent and independent forms can share the same spelling and pronunciation, as in \textit{it's \uline{his} job} (dependent) or \textit{the job is \uline{his}} (independent). Such is the case with \textit{whose}: You can ask \textit{whose gorilla is this} or \textit{whose is it}, at least when you use \textit{whose} to ask a question.

The \textit{whose} of questions is an \textsc{interrogative pronoun}, but there's another \textit{whose}, as in \textit{that's the person whose name is on the door}. This is the {relative pronoun}. Relative \textit{whose} isn't interrogating anything; it's ``standing in for'' \textit{the person's}. And the \textit{whose} in (\ref{ex:gorilla})~-- the one that Hankamer and Postal's marked as ungrammatical~-- is a relative pronoun. Putting it all together, then, the \textit{whose} at issue is this independent genitive relative pronoun.

The point Hankamer and Postal were making was that, where most pronouns have paired dependent and independent genitives, \textit{who} feels like it's missing the independent one. All the other combinations in (\ref{ex:whose-good}) are possible. Why not (d)?

\ea \label{ex:whose-good}
 \ea[]{\textit{Whose gorilla is this?}\hfill[dependent interrogative]}
 \ex[]{\textit{Whose is this?}\hfill[independent interrogative]}
 \ex[]{\textit{someone whose gorilla had left}\hfill[dependent relative]}
 \ex[*]{\textit{someone whose had left}\hfill[independent relative]}
 \z
\z

They didn't know, and that's why they published ``Whose gorilla'' as a squib. Generally, squibs don't get cited much, and this one, true to form, has been mostly ignored. Nevertheless, in 2002 the \textit{Cambridge grammar of the English language} agreed with its judgment, marking (\ref{ex:CGEL-whose}) with an asterisk. And, nearly 20 years after that, almost 30 since the article was published, \citet{Cinque2020a} was still referring to the issue a ``residual puzzle''. 

\ea[*]{\textit{The police are trying to contact the person \ob whose it was\cb.} \hfill(p. 472)}\label{ex:CGEL-whose}
\z

\bigskip

\begin{quote}
    



Linguistic discoveries often have unexpected origins. This particular journey began with a conversation about gorillas, meandered through book reviews and academic papers, and culminated in a cross-linguistic insight about the nature of possessive pronouns. Here's how it played out.

Some years ago, my friend, co-author, and mentor Geoff Pullum told me about an intriguing paper from 1973. Two linguists, Jorge Hankamer and Paul Postal, had made a bold claim in a short paper called a ``squib'': English, they said, doesn't include a particular kind of \textit{whose}. This peculiar assertion stuck in my mind, resurfacing last summer as I reviewed Geoff's most recent book \textit{The truth about English grammar} \citep{pullum2024truth}, in which he mentioned this missing \textit{whose}.

As I embarked on writing this book, I found myself drawn to linguistic edge cases~-- those fascinating exceptions and rarities that push at the boundaries of our understanding of language. The case of the independent relative \textit{whose} seemed like a perfect candidate for investigation.

The plot thickened earlier this year when I read a paper by Adele Goldberg and her graduate student Nicole Cuneo \citep{Cuneo2023} about island extraction in relative clauses (See \ref{sec:extraction-islands}). Their work sparked my interest in the role of pragmatics in these linguistic puzzles. It occurred to me that similar forces might be at work on the missing \textit{whose}, so I suggested to Adele that we team up on trying to figure this out. Although a collaboration didn't materialize due to Adele's existing commitments, I decided to pursue this line of inquiry on my own.

What followed was an intense two-week period of research and writing, during which several key insights~-- and a paper \citep{reynolds2025whose}~-- emerged, sparked in a large part by a few key observations I found in \citet{Halliday1976cohesion} and \citet{Ariel2014}.

As the pieces fell into place, I realized that my findings might have implications beyond English. A preliminary investigation into French, German, Spanish, Persian, and Japanese indeed suggested that similar constraints might be at play across languages. But before we explore these broader implications, let's go back to the beginning: Hankamer and Postal's bold claim about a supposedly non-existent \textit{whose} in English.
\end{quote}

How had Hankamer and Postal determined that (\ref{ex:gorilla}) was ungrammatical? Almost certainly by noticing their own metacognitive feelings. Did they search for instances or present relevant examples to a representative sample of English speakers to see how the independent \textit{whose} was perceived? They don't say, but it's very unlikely. Such experimental methods were almost never seen as necessary at the time. A linguist's personal  grammaticality judgments were generally deemed valid and reliable. Rarely were they questioned.

It occurred to me, though, to do something they could not have done. In March of 2024, I asked OpenAI's GPT-4 and Google's Gemini Advanced to rate the grammaticality of this \textit{whose} on a scale of 1--7, not because I questioned the feelings of Hankamer and Postal and the Cambridge grammar authors, but because I was curious to see how these large language models would respond. The sentence I gave them was (\ref{ex:whose-wasn't}).

\ea[]{\textit{Mine was working, but I know someone whose wasn't.}} \label{ex:whose-wasn't}
\z

Surprisingly to me, both models rated it a 6, denoting high grammaticality, and cited the slightly unusual use of \textit{whose} as the reason for not awarding a 7. I'd experimented with eliciting grammaticality judgments from large language models before with various construction, and I'd mostly found their assessments to align surprisingly well with my own. I'd even asked them to produce examples at various points on the scale, and again, they'd fared reasonably well. This is why I was rather surprised at a 6 rating for (\ref{ex:whose-wasn't}).

Intrigued by this mismatch, it occurred to me to ask some non-linguist colleagues and family members what they thought of it. I sent out a few quick texts and within a day, ten of the twelve I'd asked had also given it a 6 or 7. Two gave it a 4. Many didn't mention \textit{whose}, and even after I queried the pronoun specifically, two said it seemed fine to them. To my surprise, the large language models were a better for with the naive humans than were the linguists. I was intrigued.

This issue arose while I was working on \textit{Wikipedia}'s \href{https://en.wikipedia.org/wiki/English_pronouns}{English pronouns} page. There, the other editors and I had reflected the position of Hankamer and Postal and the Cambridge grammar: that independent relative \textit{whose} did not exist. With these new grammaticality judgments, I started to wonder if I should revise that entry.

\textit{Wikipedia} requires sources, though, and ``Brett's friends and family along with some LLMs'' doesn't cut it. It was then that it occurred to me to consult \textit{The Oxford English Dictionary}. Had Hankamer and Postal done so, perhaps they would not have published because they would have found that the lexicographers at the \textit{OED} seem to take the same position as my informants and the large language models: the independent genitive relative \textit{whose} is rare, but grammatical. The \textit{OED} presents examples going back to the 1300s, including those in (\ref{ex:OED-whose}). %even the first edition had these examples.

\ea \label{ex:OED-whose}
 \ea[]{\textit{After þe wille of him \uline{hos} þe werkes bez.} `According to the will of him whose works are.'\hfill[1325]}
 \ex[]{\textit{The man \uline{whose} these are.}\hfill[1611]}
 \ex[]{\textit{The most generous patron \dots~continues to be our \dots~Pittsburg and Allegheny auxiliary, \uline{whose} also is the largest contribution to the rent of our Association's new office.}\hfill[1904]}
 \ex[]{\textit{Everything depends on the person \uline{whose} this administration is.}\hfill[2018]}
 \z
\z

\noindent And (\ref{ex:online-whose}) presents a few more I found online.

\ea \label{ex:online-whose}
 \ea[]{\textit{In California you usually include a copy of the check with the offer. Earnest money despoit (EMD) can be anywhere from 1/2\%-3\% of the sales price. Mine was 3\%, I just helped a person \uline{whose} was 2\%, I've never seen an EMD less than 1/2\% in the past 3 years in California (from Shasta Co. to Orange Co.).}}
 \ex[]{\textit{I have met a person whose SSN was only 10 away from mine (yes only the second to last digit was one higher than mine) and he was born in a different state and was at least 7-10 years older than I was. And I met another person \uline{whose} was kinda close to mine but again we had nothing in common pertaining to our birthday or location of birth/SSN registries.}}
 \ex[]{\textit{Mine is also submitted on June 30 and no updates yet, i know a person \uline{whose} was filed on June 30 as well in Regular mode and he got approval on 8/16. I think we can expect a decision in the next few weeks.}}
 \ex[]{\textit{I'll be scouring eBay for an Alpha keypad. Debating whether to swap it out for the existing one or just use it for programming. The control panel isn't as hard to get to as the person \uline{whose} was installed in the top of a closet.}}
 \ex[]{\textit{Head shaved? Mine wasn't but the guy \uline{whose} was had to give a pubic sample.}}
 \ex[]{\textit{In my opinion, my 50\%/50\% bud/trim mix came in third place. One guy, \uline{whose} was the best in my opinion, he treated his trim exactly like he did his buds.}}
 \ex[]{\textit{When I entered in 1977, I actually had the longest ash, but mine was bent. I came second to a woman \uline{whose} was straight.}}
 \ex[]{\textit{I have vasovagal syncope (a specific thing triggers fainting) and it's not uncommon for a trigger to be urination or defecation. My trigger is medical stuff like needles, which mostly people are sympathetic to. But I did know another woman \uline{whose} was set off by morning urination.}}
 \ex[]{\textit{Meanwhile the original person \uline{whose} this photograph is\dots}}
 \z
\z

Based on this evidence, I revised the \textit{Wikipedia} article to include the two views. But, of the two positions, was it mine whose was right?

This is the problem with our feelings of (un)grammaticality. Hankamer and Postal are top-notch syntacticians, legendary even. And John Payne and Rodney Huddleston, the authors of the chapter of the \textit{Cambridge grammar of the English language} that marks (\ref{ex:CGEL-whose}) as ungrammatical, are among the best English grammarians the world.

On the other hand, there's no theoretical reason to expect relative \textit{whose} to go AWOL, there's the data from the OED and my web searches, there's the intuition of my friends and family, and there's whatever ineffable average over trillions of words of English is represented in the judgments of the large language models. 

So, we have two seemingly contradictory forces at play. On the one hand, certain established linguistic authorities tell us that they feel the sentence \textit{The person whose wasn't working} is ungrammatical. On the other hand, the feelings of various people, lexicographers and non-linguists alike, backed up by the vast echo chamber of everyday speech reflected in online examples and the judgments of generative AI models, suggest it might be fine.

\bigskip

Linguistics, like any field of study, bears witness to scholarly disputes. Yet, let truth be told, our squabbles over the fringe cases of English grammar~-- the quirky pronoun, the oddly placed phrase~-- are generally far removed from our everyday experience of language. Situations in which the grammaticality of the relative independent genitive pronoun might even be at issue, for instance, are vanishing rare.

Unless I've messed up, every sentence I've composed in this book is uncontroversially grammatical. The same is true of almost everything you hear and read every day and all that you say, at least in a language you grew up using or are proficient in. Still, there are edge cases like this \textit{whose} that aren't so clear cut, and others that aren't even so rare. What are we to do with them? How are we to adjudicate the situation when grammaticality is disputed, individually, as a community of language users, or as linguists and grammarians? What do we do when our feelings don't match up? Geoff Pullum has a suggestions.

Pullum was ``the worst sort of lazy and undisciplined teenager,'' eventually dropping out of high school. After a short stint as a shop assistant in a London bookstore, be joined a travelling band as the piano player. Somehow, this led to his appearing on two of the biggest selling UK albums of the 1960s. Despite his change in fortunes, he found touring tedious. The band broke up, and he turned his focus on linguistics, where his intellect, broad interests, and workaholic tendencies could be fully engaged. He began writing for the public in the early 2000s on the \textit{Language log} blog, which, at its peak, exceeded a million page views per month in part because Pullum is one of the best and funniest writers in the discipline. On top of that, he's also a philosopher of linguistics, and his philosophical contributions to the field are noteworthy.

\citet{Pullum2019a} suggests we employ \textsc{reflective equilibrium}. That is, he says that we should carefully consider our initial feelings about the grammaticality of a sentence, and then weigh these against established grammatical principles and the judgments of other competent speakers, going back and forth between these sources of information until our position seems to have arrived at a kind of resting place. If our feelings conflict with these principles or with the consensus of other speakers, we should be willing to revise our judgments. At the same time, if a grammatical principle seems to contradict the strong feelings of many native speakers, it may be the principle that needs to be revised.

The idea of reflective equilibrium, along with the term itself, is often credited to John Rawls, who advocated it as a method for establishing principles of justice. But Rawls was building on the ideas of Nelson Goodman, who was more interested in science. Specifically, Goodman was trying to understand how we could justify deductive logic.

Pullum realized that the study of grammar by linguists involves a unique blend of these two perspectives. On one hand, like scientists and logicians, linguists must deduce grammatical patterns from actual data~-- in this case, the language produced by speakers, despite the presence of errors. On the other, similar to ethicists and political philosophers, they must also consider our feelings about how language ought to be used, while recognizing that our feelings of ungrammaticality may be shaped by idiosyncratic factors like dialect, education, or even momentary inattention. Navigating the tension between these two aspects requires a process of constantly refining both the observed patterns and the guiding principles until they align in a satisfactory manner.

Applying reflective equilibrium to the case of the relative independent genitive pronoun \textit{whose}, we might proceed as follows: First, we note, as Hankamer and Postal did, our own feelings about a noun phrase like \textit{the person whose wasn't working}. For some people, as we've seen, this will feel ungrammatical. We then consider this evaluation in light of the general principles and patterns of English grammar. We might wish to start by assuming prefect regularity within paradigms, but we would quickly be disabused of this idea.

Famously, plural and past-tenses forms can be highly irregular. Less commonly recognized is that \textit{be} is the only verb with two past-tense forms: \textit{was} and \textit{were} and that a few verbs are defective: \textit{beware} for example. Among the modal auxiliary verbs, there are present- and past-tense forms for \textit{will}/\textit{would}, \textit{can}/\textit{could}, \textit{shall}/\textit{should}, and \textit{may}/\textit{might}, but no such counterpart for \textit{must}. The pronoun \textit{--self} forms are typically made up of the dependent genitive + \textit{self}, as in \textit{myself}, \textit{yourself}, \textit{herself}, and \textit{ourselves}. But when it comes to \textit{himself} and \textit{themselves}, it's \textit{self} together with the accusative form. The examples could be multiplied, but suffice it to say that rejecting an idea simply because it involves divergence from a pattern is a non-started. Still, it's a piece of the puzzle.

We would also need to consider the judgments of other speakers. If we find that a significant number of people find the sentence acceptable, we might need to question whether our initial intuition was correct, or whether the grammatical principle needs to be modified. We might also consider whether there are other constructions in English that seem to violate the principle in a similar way, and if so, how we judge their grammaticality.

Through this process of weighing intuitions, principles, and usage data, we can arrive at a more considered judgment. We may conclude that the sentence is indeed ungrammatical, and that our initial intuition was correct. Or we may decide that it represents a rare but acceptable variant, perhaps limited to certain dialects or registers. Or we may even conclude that the grammatical description we've relied on needs to be revised to accommodate this and similar constructions.

The key point is that reflective equilibrium provides a rational procedure for adjudicating disputes about grammaticality. It is not a mechanical algorithm that always yields a clear ``yes'' or ``no'' answer, but rather a method for navigating the complex interplay of intuition, systematic principles, and actual usage that characterizes our linguistic competence. By engaging in this process, both individually and collectively, we can arrive at grammaticality judgments that are not merely arbitrary opinions, but well-considered reflections of the deep patterns of our language.

\bigskip

Recall that Sampson's view that all is grammatical and Chomsky's early generative work, which assumes that something is grammatical if, and only if, it can be produced by the grammar. As distinct as they seem, these two ideas can united. All we need to do is to create a new rule every time we observe a new pattern. If we don't like Sampson's view, then we need something to constrain us, but it can't be whether or not something is generated by the grammar because that's what we're trying to decide.

in the context of ethical reasoning, serves as a method for reconciling our intuitions with established principles. In ethics, it involves adjusting our moral intuitions and theoretical principles until they cohere in a mutually supportive framework. Translated into the realm of linguistics, this means aligning our intuitive sense of what sounds right in language with the grammatical rules and theories developed by linguists. This process is not static; it requires ongoing adjustment as new linguistic data emerges and our language intuitions evolve. By adopting this philosophical approach, we extend its application from moral philosophy to the dynamic and often subjective domain of language, offering a systematic way to navigate the complexities of grammatical judgment.

Picture being in a maze. Each choice you make is a bet between what feels right and what's traditionally accepted as right. The grammar guidebook tells you there's only one way through this maze. But what if this guide is outdated, or the maze itself shifts and changes, just like language does? Goodman's idea of reflective equilibrium is like having a compass in this scenario, guiding us not by the old map, but by a constant conversation between our gut feelings about language and the grammatical structures we've learned.

With the whose conundrum as our example, it's like starting off thinking, \textit{The person whose wasn't working sounds alright to me.} But then, you bump into the expert opinions, the grammar guides, maybe even a linguistics paper or two that disagree. It's not about tossing your initial gut feeling aside; it's more about weighing it against these authoritative voices. This back-and-forth, this dialogue between what we instinctively feel and what we're told, is at the heart of reflective equilibrium. It might lead us to adjust our views a bit, maybe deciding the sentence is a tad off but not a grammatical catastrophe. Or, we might stick to our guns, now more aware of the broader conversation around it.

Yet, here's the thing: language is always on the move. It's alive, constantly evolving. New slang makes its debut, meanings drift, and yesterday's grammatical faux pas can become today's hot new phrase. Reflective equilibrium gives us a method to navigate this evolution, not with strict rules but with an open mind and a critical eye.

Consider the case of split infinitives, like the infamous \textit{to boldly go}. Grammar purists might balk, but there's an argument to be made for such constructions if they enhance the sentence. Reflective equilibrium encourages us to balance traditional rules against the effectiveness and clarity of expression in specific contexts.

Reflective equilibrium, then, isn't about finding a one-size-fits-all answer. It's about understanding the dynamics behind our language choices and settling on what sits right with us, all while acknowledging that language is a constantly changing landscape.

As we move forward, we'll take a closer look at this balancing act, drawing on real-life examples to illustrate how reflective equilibrium can help us navigate the complex terrain of grammatical judgment. It's about getting comfortable with linguistic ambiguity and appreciating language's rich, ever-changing nature.

\section{Free adjuncts, danglers, and howlers}

Dangling participles haunt the dreams of grammarians. \textit{\uline{Turning his attention back to his work}, a dangler stared him in the face.} As you first read this, you develop the idea that somebody was turning his attention back to his work, and, right after the comma you see \textit{a dangler}, and in your developing sense of the meaning, you entertain the idea that that dangler was that person. And then you encounter \textit{him}, which forces a complete reassessment of the situation. This is the experience of reading a bad dangling modifier, or what linguists call \textsc{free adjuncts}.

Consider (\ref{ex:erection}), noticed by my friend, Jim \citet{Donaldson2020}.

\ea[*]{\textit{Suddenly, he sits up and tugs my panties off and throws them on the floor. Pulling off his boxer briefs, his erection springs free. Holy cow!} \\ \hfill(E. L. James \textit{50 Shades of Grey})}\label{ex:erection}
\z

This dangler is bad. It qualifies as a howler, and, in response, I've marked it with an asterisk. I think it's ungrammatical, but is it?

The first example cited by The Penguin Dictionary of American English Usage and Style (by Paul W. Lovinger; New York: Penguin Reference, 2000) is as clear a case as one could want, taken from a book about cannabis (if you could steel yourself, please, I want no politically incorrect giggling at this one):

\ea
 \ea[]{\textit{Although widely used by the men, Bashilange women were rarely allowed to smoke cannabis.}}
 \ex[]{\textit{Rich and creamy, your guests will never guess that this pie is light.}}
 \z
\z

\begin{quote}
``Does this fall under the no dangling modifier prescription?'', asks Rosanne, in a post on The X-bar. Yes, Rosanne, it does. Like participles, adjectives and also some idomatic preposition phrases, when used as adjuncts, need an understood subject (or, it might be better to say, a target of predication) to be filled in if they are to be understood. The prescriptive tradition says that the subject filled in must be the one obtained from the subject of the matrix clause. Here that would be your guests, which makes a nonsense reading, so the sentence cited would be treated as an error.

But the prescriptivists have a problem. Sentences of this kind, which call for you to fill in the understood subject from somewhere else (here, the subject of is light in the subordinate clause), are so common that when I and several friends have spent some time picking new ones up from print and radio sources, we get them at a rate of as many as one per day. That's in edited sources, where grammatical errors have almost entirely been screened out. This just cannot be syntactic error. It's too frequent.\hfill\citep{LanguageLog218}
\end{quote}

Donaldson's PhD dissertation, supervised by Geoff Pullum, delves into the intricacies of dangling modifiers, providing a nuanced perspective on their usage. \citep{LanguageLog218,LanguageLog1938,LanguageLog1973,LanguageLog2153} and Donaldson argue that while danglers can be seen as ``ill-written and discourteous'', ``inconsiderate'', ``stunningly inept'', or ``bad grammatical manners'', they do not necessarily constitute ungrammaticality. This stance challenges traditional views, which often categorize such constructions as clear violations of grammatical norms.


The construction means that his erection pulled off its boxer briefs, but that is obviously not the intended meaning. ``Just about every rhetoric, grammar, and handbook written since the latter part of the 19th century warns the student against such constructions'' \citep[314]{Gilman1989}.

Consider the sentence: \textit{Having arrived late, the meeting was already over.} Traditional grammar would flag this as a dangling modifier since the subject of the participle phrase \textit{Having arrived late} does not logically correspond to the subject of the main clause, \textit{the meeting}. Pullum and Donaldson's perspective invites us to reconsider whether this incongruence automatically renders the sentence ungrammatical or merely less clear or elegant.

Their argument suggests that the judgment of grammaticality involves more than a binary distinction between right and wrong. It encompasses considerations of clarity, elegance, and the speaker or writer's intent. From this viewpoint, the use of a dangling modifier might reflect stylistic choices or conversational shortcuts rather than a lack of grammatical competence.

Reflective equilibrium offers a valuable framework for navigating this debate. It encourages us to balance our intuitive disapproval of danglers with the understanding that language flexibility often accommodates such constructions in everyday usage. This balancing act involves weighing the clarity and communicative effectiveness of language against the descriptive realities of how people speak and write.

Moreover, examining the acceptance of dangling modifiers in different contexts and registers can shed light on the fluid boundaries of grammaticality. For instance, creative writing and informal speech frequently employ danglers to convey a sense of immediacy or emotional resonance, prioritizing expressive impact over strict adherence to grammatical rules.

\subsection{Singular \textit{they}}

My daughter, Mia, had two roommates, one of whom, I'll call her Sasha, wanted others to refer to her using the pronoun \{\textit{they}\}. This is easy and natural for Mia. It's fully grammatical for her. She can say \textit{Jessica cleaned up her mess but Sasha didn't. I wish they would,} meaning not that she wishes both roommates would clean it up but that Sasha would.

That's not what that means for me though. When I hear this, I understanding that Mia wants both of her roommates to clean.

Or consider this examples that I heard on \textit{This American Life}:

\begin{quote}
Things for Malinda and the kids are pretty good these days. Stella just turned 18. She says her climate nightmares only come every couple weeks instead of nightly. And she's off at college this fall, a milestone that, for a while, she never thought she'd reach. Emrys is 19. They work on a boat, teaching kids about marine life and maritime skills, a job they love, though they skip out on teaching the lessons about climate change. They said they're not ready to do that yet. \hfill\citep{thisamericanlife_748} 
\end{quote}

This understanding is ``ballistic'' for me, in the sense that it can't be stopped, at least not at the moment. It's like seeing $\langle$cat$\rangle$ and having the word \textit{cat} pop into my mind. There's nothing I can do do stop it.

I can learn a new rule for \textit{they}, but it's a clunky and artificial kludge slapped on top of what is an exquisitely tuned system, one that can do no better than to try to funnel the champagne into a different bottle as it foams forth. It's the effortful correction you make when you see \textit{the colour of the following word is} {\color{blue}\textit{red}}.

This is not because \{\textit{they}\} cannot be singular. I have no trouble at all with singular \{\textit{they}\}, as in 
\textit{Someone forgot \uline{their} bag}, but it doesn't compute for me when it refers to a named person.

My theoretical understanding of singular \{\textit{they}\} tells me it should be grammatical. If this is confusing to you, recall that singular \{\textit{you}\} is perfectly grammatical, even though \{\textit{you}\}, like \{\textit{they}\} started life as a plural pronoun. 

The singular second-person pronoun that \{\textit{you}\} displaced was \{\textit{thou}\}. As in many Western European languages, \{\textit{thou}\} didn't simply mean `the individual I am addressing.' It also conveyed a sense of familiarity even intimacy, or where that was absent, a sense of superiority of the speaker over the spoken to. It's complement, the second-person plural had the additional social meaning of respectful address.

Over time, the `familiar/intimate' meaning was overpowered by the implication that the person being \textit{thou}ed was being spoken down to. And, so, in the name of politeness, \{\textit{thou}\} was replaced with plural \{\textit{you}\}.

\{\textit{You}\} is still syntactically plural. The verb agreement is the plural agreement. But semantically, it can be singular or plural.

\{\textit{They}\} is the same. There's no reason that it could not follow the precedent set by \textit{you}. And yet, I can't get there. I keep getting blindsided by cases like the messy roommate(s).

I've tried to reach a reflexive equilibrium, but updating my feeling of grammaticality for referring to a named individual, such as ``\textit{Sasha}'', as \{\textit{they}\} isn't as easy as learning the meaning of a new word. It requires unlearning half a century of practice. And that's why we have gradient grammaticality in a population.

In trying to achieve a reflexive equilibrium, I did realize that singular \{\textit{you}\} does differ from the new singular \{\textit{they}\} in one somewhat related aspect. \{\textit{You}\} is not interchangeable with a person's name, not without infantilizing them. If I'm speaking to my daughter, I can ask \textit{what do you think,} but I can't ask \textit{what does Mia think} without a great deal of implied condescension.
\newpage
\subsection{The double-\textit{is} construction}

There's a curious construction that you hear all over the place these days. In fact, I heard it this morning on the popular \textit{EconTalk} podcast, hosted by economist Russ Roberts, and not for the first time. Roberts and his guests~-- people like Kevin Kelly, co-founder of \textit{Wired} magazine, and Eric Hoel, a research scientist and author known for his work on consciousness and theoretical neuroscience~-- will say things like:

\ea
 \ea[]{\textit{But, the problem \uline{is, is} that there's not just, say, the coasts versus the middle or the South versus the rest.}}
 \ex[]{\textit{But, my point \uline{is, is} that, in America the bureaucratic infrastructure is not well-respected.}}
 \ex[]{\textit{But, the issue \uline{is, is} that once the mask is on, it's very unclear.}}
 \z
\z

This construction, with its doubled \textit{is}, has a fancy name: the double-\textit{is} or ISIS construction. It's become surprisingly common in the last 40 years or so. Despite this, from a strictly structural viewpoint, it's a bit of an oddity.

At first glance, the ISIS construction might look like other common structures. For example:

\ea 
 \ea[]{\textit{The question \uline{is} is it dangerous?}}\label{ex:the-question-is-is}
 \ex[]{\textit{What it is \uline{is} a tool for learning about language.}}\label{ex:what-it-is-is}
 \z
\z

But there's a key difference. In these examples, each \textit{is} has its own subject and serves a clear purpose within its own clause. In (\ref{ex:the-question-is-is}), the main verb of the sentence is the first \textit{is}~-- the underlined one. In (\ref{ex:what-it-is-is}), it's the second.

With the ISIS construction, though, there are more \textit{is}es than clauses, and no-one can say with certainly which is the main verb.

\bigskip

So, why does this strange construction exist? One theory points to the focusing function of the constructions in (\ref{ex:the-question-is-is}) and (\ref{ex:what-it-is-is}) and suggests it has been misattributed to the coincidentally-shared \textit{is~is}.

What I mean by ``focusing function'' is that the examples seem a little strident, a little assertive. (\ref{ex:the-question-is-is}) doesn't just ask ``is it dangerous?'' It focuses the listener on this question, asserting that this is \textit{the} question. Similarly, (\ref{ex:what-it-is-is}) isn't just saying ``this is a tool for learning about language.'' It's focusing on that aspect: That's what the thing is.

Each example focuses our attention in its own way, but what they share is \textit{is~is}. In each case, it's like a brief dramatic pause before the main event. The switching on of the spotlights before they converge on the key message.

Compare the simple statement \textit{It's a fabulous thing} with the ISIS example \textit{The point \uline{is, is} that it's a fabulous thing.} If it were a question, it could be \textit{The question is \ob is it a fabulous thing\cb}, but it's not a question~-- it's a point. And so \textit{point} replaces \textit{question} and out pops ISIS.

\bigskip

The ISIS construction is another example of conflicting signals of grammaticality. Roberts and his guests~-- polished writers and speakers all~-- use it relatively frequently, and it mostly goes unnoticed. Others, like me, though, notice it all the time. My grammatical intuitions may be contaminated, but even many years before I knew anything about the construction, I had noticed it in a friend's speech and found it ungrammatical. Formally speaking, a second \textit{is} without a subject to connect to isn't found elsewhere in English. Also, ISIS doesn't turn up in edited writing.

I understand that it seems to be grammatical for many people, at least in spoken English. I also understand that there's no need for spoken English to conform to the grammatical norms of written English. I even know of constructions that I find grammatical despite their being unlike anything else in English, \textit{I'\uline{d better} try again}, \textit{\uline{many's} the time\dots}, \textit{\uline{what with} the way it went\dots}, \textit{\uline{no matter} the change\dots}, etc. And so, I don't think the extra \textit{is} necessarily disqualifies it from being grammatical. In fact, it is my considered judgment that ISIS must be grammatical for people like Roberts and his guests.

And yet I never find myself saying it, and it always throws up a red flag when I hear it. I've tried over many years to overcome my feeling of ungrammaticality, but it seems I'm stuck with it.

%\subsection{\textit{One less} vs \textit{one fewer}}

%The Corpus of Contemporary American English has 139 instances of \textit{one fewer} and 1,247 instances of \textit{one less}. \textit{I have one less problem to deal with} sounds much better to me than \textit{I have one fewer problem to deal with,} and yet the direction of the preference~-- if not the magnitude~-- is reversed for \textit{fewer}/\textit{less problems}: 171:45.

%\subsection{Double \textit{is} (ISIS)}

%For about the last 40 years, linguists have been noticing a particular construction. It's known as the ISIS construction or the double-\textit{is} construction and is characterized by an excess of \textit{is}es without corresponding subjects, as in \textit{The thing is is that it's kind of strange}. It has become increasingly common, especially with subjects like \textit{the thing}, \textit{the problem}, \textit{the fact}, \textit{the point}, and \textit{the truth}. Personally, I first noticed it in the speech of an American friend while working in Tokyo around 1998, before I knew it was a something people thought about. I remember feeling it was ungrammatical and pointing it out to him, but he was unmoved and still uses it today. So, is it grammatical or is it not?

%Outside of ISIS, it's rare for \textit{is} to follow \textit{is} directly~-- rare, but not impossible. What it \uline{is~is} a clause as the subject. The underlined string in \textit{I know \uline{what it is}} is a subordinate clause. Here it's a complement in the verb phrase, but it can also be the subject as in (\ref{ex:what-it-is-is}), and, although it's most common with \textit{is~is}, you can switch out either verb and the sentence still works.

%\ea \label{ex:what-it-is-is}
% \ea[]{\textit{\uline{What it is} is a combination of form and meaning.}}
% \ex[]{\textit{\uline{What I think} is that it's not important.}}
% \ex[]{\textit{\uline{Who she is} matters.}}
% \ex[]{\textit{\uline{Where it goes} makes no difference.}}
% \z
%\z

%These are perfectly grammatical English sentences. But, as you can see, when the subject of the sentence is a clause, two tensed verbs end up back to back, the first one belonging to the subject and the second to the main clause. We often find multiple verbs one after the other, as in \textit{He \uline{likes swimming}}, but rarely is the second verb tensed. That's why it's a little unusual that all the verbs in (\ref{ex:what-it-is-is}) are tensed.

%Importantly, each of the examples in (\ref{ex:what-it-is-is}) seems a little strident, a little assertive. They've been repackaged with a pumped up subject~-- a clause~-- and a correspondingly flat predicate. But these puffed up subjects can be punctured and the sentences smoothed out so that they don't draw as much attention to themselves.

%\ea \label{ex:what-it-is-is}
% \ea[]{\textit{It is a combination of form and meaning.}}
% \ex[]{\textit{I think it's not important.}}
% \ex[]{\textit{It matters who she is.}}
% \ex[]{\textit{It makes no difference where it goes.}}
% \z
%\z

%There's another construction that looks a little like this on the surface and has a similar function but has a different underlying form: \textit{The question \uline{is~is} it the link to ISIS.}

%First, just like the examples above, it has \textit{is~is}. But the similarity is merely co-incidental. In this case the subject is just \textit{the question}. The first \textit{is} above was part of the subject, but this time it's just the main verb of the clause. The second \textit{is} belongs to the question, \textit{Is it the link to ISIS?}

%Second, again, like the examples above this is a construction is a repackaging that serves to frame and highlight the upcoming information as the central focus of the sentence. They act as a signal to the listener/reader that the most important content is about to be delivered.

%These similarities lead to a question: Could the older constructions have inspired the newer one? Perhaps, people started associating strong topic--comment meaning to the \textit{is~is} string instead of to the clause as subject. And so they started doubling the \textit{is}. 

%Tensed verbs in English have subjects. Even if we've ``dropped'' them as in (\textit{it}) \textit{\uline{gets} me every time}, they're easily retrievable. If two tensed verbs are coordinated with \textit{and}, they can share a subject, as in \textit{She \uline{went} and \uline{did} it again,} but what's to be done with the extra \textit{is}es in 

%\ea \label{ex:ISIS}
% \ea[]{\textit{The thing is is that I wanted to do well.}}
% \ex[]{\textit{The fact is is that we lost.}}
% \ex[]{\textit{The fact is is that we lost.}}
% \z
%\z

\section{What does (un)grammaticality feel like 1?}

Sensations arise from the interaction between stimuli, our sense organs, and our nervous system. We've got the traditional five senses: sight, taste, smell, hearing, and touch. But it goes well beyond that. We have a sense of balance, regulated by nerves in our vestibular system in our inner ears. Our sense of hunger is triggered by nerves in the gastrointestinal tract and transmitted through the vagus nerve to the hypothalamus. We sense the location of our various body parts and the angles of our joints, pain, pressure, itch, hot and cold, and chemical irritants like capsicum and allyl isothiocyanate in wasabi and mustard. We even have a sense of the levels of carbon dioxide and oxygen in our blood.

We're sensitive creatures, but we have no sense organs for grammar, so grammar is not something we sense, and (un)grammaticality is not a sensation. But it is a feeling, or at least part of it is. Feelings are sensations that have a evaluative aspect.

\begin{quote}
    No sooner had the warm liquid mixed with the crumbs touched my palate than a shudder ran through me and I stopped, intent upon the extraordinary thing that was happening to me. An exquisite pleasure had invaded my senses, something isolated, detached, with no suggestion of its origin. And at once the vicissitudes of life had become indifferent to me, its disasters innocuous, its brevity illusory – this new sensation having had on me the effect which love has of filling me with a precious essence; or rather this essence was not in me it was me. ... Whence did it come? What did it mean? How could I seize and apprehend it? ... And suddenly the memory revealed itself. \\\hfill~-- Marcel Proust, \textit{In Search of Lost Time}
\end{quote}

Proust describes a feeling triggered by but unconnected to any immediate sensation, a feeling about a memory but without, at least initially, access to the memory itself. It is a feeling about a feeling, a metacognitive feeling.

A toddler just learning to talk has feelings about the speech of others.

``He knew only that his child was his warrant. He said: If he is not the word of God, God never spoke.''

Faulkner writes, ``Memory believes before knowing remembers. Believes longer than recollects, longer than knowing even wonders.''


``I can't go on. I'll go on.''


``In the end the Party would announce that two and two made five, and you would have to believe it. It was inevitable that they should make that claim sooner or later: the logic of their position demanded it.''

``And thus the native hue of resolution
Is sicklied o'er with the pale cast of thought,
And enterprises of great pith and moment
With this regard their currents turn awry,
And lose the name of action.''

``The heart has its reasons of which reason knows nothing.''
 ~--  Blaise Pascal


Nick Chater, in The Mind is Flat, argues that mental processes are fundamentally superficial, challenging the deep and complex models often posited in psychology and cognitive science. From this perspective, metacognitive feelings, including the sensation of (un)grammaticality, are immediate and surface-level responses, not rooted in any deep, structured understanding of grammar or cognitive processing.

Chater would likely suggest that the feeling of (un)grammaticality arises not from an intricate, rule-based analysis of language structures internalized through extensive learning, but from heuristic responses generated by the brain's superficial processing systems. These systems are quick to judge based on patterns and examples encountered in everyday language use, rather than deep grammatical rules. For instance, a sentence might feel ungrammatical simply because it's unfamiliar or deviates from frequently observed patterns, not necessarily because it violates a specific grammatical rule.

This approach emphasizes how our perceptions of grammaticality are shaped more by pragmatic cues and less by any robust internal grammar engine. Thus, the feeling of (un)grammaticality is a product of our minds navigating through language based on surface impressions and immediate judgments, aligning well with Chater's view that the mind constructs its seemingly complex operations from moment to moment, without any underlying depth.

\section{What does ungrammaticality feel like2?}

Encountering an ungrammatical sentence can feel like hitting a speed bump while cruising down a linguistic highway. It's a jolt, a disruption, a sudden break in the flow of comprehension. Your brain, which was humming along smoothly, parsing predicates and conjoining clauses, suddenly grinds its gears, trying to make sense of the syntactic wreckage before it.

But what, exactly, \textit{is} this feeling of ungrammaticality? Is it a sharp pain, like the one you get when you stub your toe on the irregular past tense of \textit{lead}? Is it a queasy unease, the grammatical equivalent of carsickness, brought on by a winding road of nested subordinate clauses? Or is it more like the feeling of stepping off a curb you didn't expect, a momentary loss of balance as your linguistic footing shifts?

Psycholinguists and neuroscientists are just beginning to grapple with these questions, to probe the subjective experience of processing language that deviates from internalized norms. And what they're finding suggests that the feeling of ungrammaticality is not one thing, but many.

Consider the sentences in \ref{ex:green-cheese}. They're all ungrammatical, but they're likely to elicit different reactions. The first might provoke a double-take, a moment of confusion as your brain tries to untangle the semantic anomalies. The second might feel more jarringly wrong, the linguistic equivalent of nails on a chalkboard. And the third? It might just make you laugh, a ``word salad'' so thoroughly scrambled that it short-circuits your brain's error-handling mechanisms and tips you into the absurd.


\ea\label{ex:green-cheese}
    \ea[*]{\textit{The cheese green ideas sleep late furiously.}}
    \ex[*]{\textit{Late the ideas sleep furiously green cheese.}}
    \ex[*]{\textit{Sleep furiously late ideas the green cheese.}}
    \z
\z

These varied responses hint at the complex interplay of factors that shape our subjective experience of ungrammaticality. There's the severity of the violation, ranging from minor syntactic stumbles to major structural catastrophes. There's the level of processing at which the error occurs~-- a semantic mismatch might feel different from a morphological mistake. And then there's the role of expectation and surprise~-- an error that subverts a strongly predicted continuation might land with more of a jolt than one in a more uncertain context.

But amidst all this variability, there's a curious commonality: while ungrammatical sentences often evoke visceral reactions, \textit{grammatical} sentences rarely do. A sentence that hews flawlessly to the rules of syntax~-- as almost all we encounter do~-- don't typically elicit a feeling of rightness to match the wrongness of their ungrammatical counterparts. A positive feeling does not accompany each grammatical utterance. Grammaticality, it seems, is the linguistic baseline, the unremarkable background against which errors stand out in sharp relief. It just feels normal.

There are exceptions, of course. A particularly mellifluous line of poetry, a perfectly crafted witticism, a sentence of such crystalline clarity that it seems to distill the essence of a complex idea~-- these might provoke a feeling of aesthetic appreciation, a sense of language used masterfully. And for language sticklers and grammar mavens, a well-wrought correction might evoke a certain righteous satisfaction, a sense of order restored to a chaotic linguistic universe.

But for the most part, grammaticality is like clean air or a healthy back~-- you don't notice it until it's gone. It's the linguistic equivalent of homeostasis, the stable equilibrium that allows us to extract meaning effortlessly from the words we read and hear. Only when that equilibrium is perturbed~-- by a syntactic stumble, a semantic contradiction, or a morphological mismatch~-- do we feel the friction of processing gone awry.

So why this asymmetry between our responses to the grammatical and the ungrammatical? One possibility, suggested by recent work in the predictive processing framework \citep{Clark2013,Friston2010}, is that it reflects the brain's fundamental mode of operation~-- not as a passive recipient of sensory input, but as an active, predictive engine, constantly striving to anticipate and make sense of the world around it.

According to the predictive processing framework, the brain constantly generates top-down predictions about the likely causes of its sensory inputs, based on its prior knowledge and experience. These predictions are then compared to the actual bottom-up sensory data, and any mismatches between the two generate prediction errors \citep{Clark2013,Friston2010}.

In the domain of language, this means that the brain is continually making predictions about the likely structure and meaning of upcoming words and phrases. These predictions are based on the statistical regularities of the language it has been exposed to, as well as the current linguistic context. For example, if you hear the phrase \textit{The cat sat on the\dots}, your brain automatically predicts that the next word is likely to be something like \textit{mat}, \textit{couch}, or \textit{carpet}. This prediction is based on your prior knowledge of English phrase structure, as well as your experience with the typical things that cats sit on \citep{Kuperberg2016}.

Crucially, these linguistic predictions are being made at multiple levels simultaneously. At the lowest level, the brain is predicting the likely phonological or orthographic features of upcoming words. At higher levels, it's predicting their likely syntactic categories, semantic features, and even pragmatic implications. All these predictions are constantly being updated and refined as new linguistic input comes in.

Now, here's the key point: when the brain's predictions are correct~-- which is most of the time, given the highly predictable nature of language~-- it doesn't need to do anything. We could say the matching input ``cancels out'' the prediction, and the brain carries on more or less as if nothing had happened. In a sense, the brain is already ``ahead of the curve'', having successfully anticipated the input before it arrives. As a result, no update is needed. 

This successful prediction is a sign of a well-adapted, efficiently functioning language system. It means that the brain's internal model of the language is well-calibrated to the actual statistics of the linguistic environment. ``God's in his Heaven, All's right with the world!'' And it's this successful prediction that allows language comprehension to be such a fast, automatic process~-- we don't have to consciously work out how each word and phrase relates, because our brains have already done the work for us.

Interestingly, this absence of prediction error when things are going as expected may explain the curious phenomenon we discussed earlier~-- the fact that grammatical sentences don't usually evoke any particular feeling of ``rightness''. When a sentence perfectly matches our brain's grammatical and semantic predictions, there's no error signal to be consciously perceived. The linguistic processing is so smooth and effortless that it doesn't rise to the level of conscious awareness.

It's only when the brain's predictions are wrong~-- when the incoming language deviates from our expectations~-- that we become consciously aware of the linguistic processing. And this is precisely because these deviations generate prediction errors~-- neural signals of surprise that propagate up the processing hierarchy, triggering conscious perception and higher-level cognitive processes like reanalysis and reinterpretation \citep{Rabovsky2018}.

\bigskip

In \S\ref{sec:tristan}, I told the story of the Tristan chord, F, B, D♯, and G♯ from Wagner's opera ``Tristan und Isolde''. Audiences at the time, had no experience with such a chord. It was entirely unpredictable, and so it registered as wrongly dissonant, the musical equivalent of ungrammatical. But some where able to find a sense in it, a narrative sense, rather than a harmonic one. Later, I'll get into how constructions and other linguistic forms can mean things other than denotations and predications, things like identity, status, deference, interrogation, etc., but the point is that meaning isn't a single thing. And because some people could see analogies between unfulfilled harmonic tension and the unrequited love central to the opera's story, it wasn't dismissed.

Almost certainly it was the case that people had to work to find such a sense, and they did so because the author of the chord was Wagner. This is crucial context. Wagner had come to be seen as a musical genius, and his previous operas~-- notably Der fliegende Holländer (1843), Tannhäuser (1845), and Lohengrin (1850)~-- had featured increased chromaticism and a more flexible approach to tonality. He was also renowned for his concept of the Gesamtkunstwerk, or ``total work of art,'' which sought to unify all artistic elements within an opera. So, while the musical context of the Tristan chord evoked a sense of wrongness, this sense may have been attenuated by the context of having Wagner's name attached to it.

A lesser musician~-- someone young and unknown, or one with other such prejudicial strikes against them~-- may have been seen as simply incompetent had they composed or performed such a thing. But, although the chord didn't make sense, nor did it make sense to believe that the great Opernkomponist had made error such as this. It had to be intentional, and so the task these listeners set themselves was to determine the intention. And they may have set off on that journey because the predictive brains had dampened the feeling of dissonance for these folks because, well, ``it's Wagner''?

Whether my stylized story is true or not, what is indisputable is that the opera continued to be performed. While it was not an immediate sensation, it gradually gained recognition and acclaim for its innovative music and dramatic storytelling.

The opening run itself was not extended, but that was common for opera premieres at the time. Subsequent productions in other cities helped to build its reputation. By the 1870s and 1880s, ``Tristan und Isolde'' had become a staple of the operatic repertoire and was widely regarded as a masterpiece. And, as a result, more and more people were exposed to the chord. More composers started using it, exposing even more audiences, and folks got used to it. It came to reside in the predictable, and it came to sound no wronger than a minor chord in the right contexts.

\bigskip

So in a real sense, the feeling of ungrammaticality is the flip side of the brain's successful prediction of grammaticality. It's the conscious signal that something has gone wrong, that the brain's linguistic model needs updating. And it's this constant cycle of prediction, error, and updating that allows our brains to so efficiently navigate the complex linguistic world.

But the real beauty of this system is that even the errors are useful. Each time the brain is surprised by an ungrammatical or semantically anomalous sentence, it's an opportunity for learning. The prediction error signals the need to update the brain's linguistic model, to tweak its expectations to better match the actual statistics of the language. Over time, this constant fine-tuning leads to an increasingly sophisticated and well-calibrated predictive model, one that can efficiently handle even the most complex and ambiguous linguistic inputs.

In this way, the predictive processing framework offers a powerful and unifying account of the cognitive and neural dynamics of language~-- an account that sees linguistic processing not as a passive reception of information, but as an active, hypothesis-testing process of prediction and error-driven learning.

When expectations are met~-- when the linguistic input matches the brain's predictions~-- the result is a smooth, largely unconscious flow of language processing, unperturbed by prediction error. It feels like nothing, like baseline. But when those expectations are violated~-- when the brain encounters a word, a phrase, or a sentence that deviates from its predicted pattern~-- the result is a cascade of prediction errors, a neural signal of surprise and updating that's consciously experienced as a negative feeling of ungrammaticality \citep{FernandezVelasco2021}.

\bigskip

Interestingly, this predictive processing account also offers a way to understand the contextual variability of grammaticality judgments. If our sense of ungrammaticality arises from the brain's linguistic predictions, and if those predictions are shaped by the statistical regularities of the language environment, then it follows that individuals with different linguistic backgrounds may have slightly different intuitions about what counts as grammatical or ungrammatical.

For example, double negatives such as \textit{I don't know nothin about that} or verb forms such as \textit{I seen him yesterday} are familiar to most English speakers, but they're not part of Standard English. As a result, in a context where Standard English prevails, they may provoke a feeling of ungrammaticality. But when the Kid Rock laments, ``I ain't seen the sun shine in 3 damn days'' or Lil West boasts, ``I don't take no L's,'' we don't expect Standard English, and no feeling of ungrammaticality is evoked.

Young children over-regularize past tenses. So, instead of \textit{I found you,} a child might say, \textit{I finded you}. On the one hand, this is ungrammatical and we all know it. On the other, it doesn't feel as ungrammatical as if I said to you \textit{I wrote two book}. Even if we hadn't expected the child to say \textit{finded}, the reason for it is readily apparent in the context, and our sense of ungrammaticality is commensurately dampened. Similarly, should we hear someone with a noticeable foreign accent say, \textit{I must to go,} it registers as ungrammatical, but at the same time, the level of surprise is attenuated by the inference that this is someone who has come to English later in their life.

Conversely, I speak Japanese with a reasonable level of conversational fluency, but if Japanese speaker confidently said (\ref{ex:fond-of-cats}b) to me instead of (a) to explain why she had a cat, I probably wouldn't notice or experience any feelings of ungrammaticality. In (a), \textit{neko} `cat' is the topic~-- thus \textsc{top}~-- and the lack of a subject suggests that a first-person subject `I'. \textsc{Cop} means a copula, roughly \textit{be}, so in a stilted overly direct translation it's `as for cats, I'm fond (of them).'

\ea\label{ex:fond-of-cats}
    \ea[]{\textit{猫~~~~~} \textit{は~~~} \textit{好き} \textit{です}。\\
      \gll Neko wa suki desu. \\
           cat \textsc{top} fond \textsc{cop} \\
      \glt `I'm fond of cats.'
      }
    \ex[]{\textit{猫~~~~~} \textit{が~~~} \textit{好き} \textit{です}。\\
      \gll Neko ga suki desu. \\
           cat \textsc{subj} fond \textsc{cop} \\
      \glt `Cats are fond of it.'
      }
    \z
\z

The second sentence, though, has \textit{cat} as the subject (\textsc{subj}), with no particular topic stated. That means the topic is probably whatever we've been talking about, and the sentence tells us that cats are fond of whatever it is. I know this \textit{ga}--\textit{wa} distinction, and I usually make the right choice when speaking. But English doesn't mark topics vs subjects this way, and so it might just slip by me if I wasn't primed for it.

\bigskip

The sense of ungrammaticality we feel is not a binary, all-or-nothing affair. It's a graded experience, one that spans highs and lows, that reflect the uncertain nature of prediction. A minor syntactic slip, something very close to expectations, might generate a small, localized prediction error, experienced as a fleeting sense of discomfort. A more serious grammatical deviation might generate a larger, more widespread error signal, experienced as a visceral sense of linguistic wrongness or confusion.

And yet what will register as serious isn't easy to pin down. Subject--verb agreement errors seem objectively minor, often rectifiable through the mere toggling of a single \textit{--s}, but \uline{this error feel} egregious. Other times, such disagreements might slide by with hardly a niggle. Perhaps you sensed the infelicity in the previous paragraph: \textit{reflect} should have been \textit{reflects}. If you did, I'd guess the intensity of ungrammaticality you felt was minor compared with \textit{this error feel\dots}.

This feeling can actually be modelled quantifiably as (\ref{eq:Bayes}).

\begin{equation} \label{eq:Bayes}
P(A|B) = \frac{P(B|A) \cdot P(A)}{P(B)}
\end{equation}

This equation, known as Bayes' theorem, is a fundamental principle in probability theory and statistics. In the context of language processing, it provides a mathematical framework for understanding how the brain updates its predictions based on new linguistic input.

Here's how we can interpret the terms in the equation:

\begin{itemize}
    \item $P(A|B)$ is read ``the probability ($P$) of $A$ given $B$''. It denotes, on a scale of 0 to 1, the probability of a particular grammatical construction $A$ being used, given the current linguistic and situational context $B$. This is what the brain is trying to predict.
    \item $P(B|A)$ is read ``the probability ($P$) of $B$ given $A$'' and is the likelihood, the probability of the current context $B$ obtaining when the construction $A$ is used. This is based on our prior experience with the language~-- how often we have been in this context when construction $A$ is used.
    \item $P(A)$ is the probability of construction $A$, based on its overall frequency in the language, regardless of context.
    \item $P(B)$ is the probability of the current context occurring, regardless of which construction is used.
\end{itemize}

When the brain encounters a new linguistic input, it updates its prediction of the most likely construction based on how well that construction fits with its prior expectations and how well it explains the current context. An example might be useful.

Let's assume we're in a kitchen with a cat and a bowl on the floor with a drop of milk in the bottom, and I've said \textit{the cat}. This is the situation $B$. Your brain has already started guessing what I'll say based on what's highly probable given $B$. Two possibilities are \textit{The cat drank the milk}, which we'll call $A_1$, and \textit{the cat the milk drank}, which we'll call $A_2$. $A_1$ is a grammatical sentence that follows a typical main-clause structure in English. $A_2$ is a grammatical noun phrase could be a clause with the same meaning as $A_1$, but it's far from the typical main-clause structure. Alternatively, it could be a grammatically standard noun phrase denoting a feline in the rather odd state of having been drunk by milk.

In reality, your brain's model would not be limited to choosing between these two options, but this simplification will serve for our example.
Based on the frequency of these structures in your experience, you might have following probabilities:

\begin{enumerate}
    \item $P(A_1)$ = 1.0 (the probability of the clause)
    \item $P(A_2)$ = 0.0000000001 (the probability of the NP)
    \item $P(B)$ = 1.0 (the probability of your being in the current situation)
    \item $P(B|A_1)$ = 0.01 (the probability of being in the current situation if you hear $A_1$)
    \item $P(B|A_2)$ = 0.000000000000001 (the probability of being in the current situation if you hear $A_2$)
\end{enumerate}

The 1.0 probability of $P(A_1)$ means that you've heard this clause before. You may even have uttered it yourself. The NP, though, is unfamiliar, and its meaning makes it highly unlikely, even if its structure is perfectly normal. Your brain puts it at $10^{-10}$. Next, to make things easy, we'll say the probability that we're in the situation is 1.0 because we do indeed find ourselves in this situation. And, contextually speaking, the chance that I'll say \textit{the cat drank the milk} is nowhere near 1, but it's not that low either. On the other hand, the chance I'll say \textit{the cat the milk drank} seems tiny. You've given it a probability of $10^{-15}$.

We can now calculate the probability that I say \textit{the cat drank the milk} using equation \ref{eq:cat-drank-milk}.

\begin{equation} \label{eq:cat-drank-milk}
P(A_1|B) = \frac{P(B|A_1) \cdot P(A_1)}{P(B)} = \frac{0.01 \cdot 1.0}{1.0} = 0.01 \times 1.0 = 0.01
\end{equation}

And the probability that I say \textit{the cat the milk drank} can be calculated using equation \ref{eq:cat-milk-drank}.
\begin{equation}\label{eq:cat-milk-drank}
P(A_2|B) = \frac{P(B|A_2) \cdot P(A_2)}{P(B)} = \frac{10^{-15} \cdot 10^{-10}}{1.0} = 10^{-15} \times 10^{-10} = 10^{-25}
\end{equation}

So, given the two options, \textit{the cat drank the milk} is overwhelmingly more probable in this situation~-- by 23 orders of magnitude~-- though there is a good deal of room for me to say something else entirely, such as \textit{the cat seems to like milk}, or \textit{the cat may be planning a coup to overthrow the household's leadership}, or just \textit{the cat}. After all, cats are full of surprises and, apparently, so is Bayesian inference. Nevertheless, you're now predicting the completion \textit{drank the milk} because, even though it's not particularly probable, it is the most probable of the alternatives.

These probabilities represent your brain's expectations before I finish my sentence. Now, let's say, dullard that I am, that I do indeed blurt out, ``The cat drank the milk.'' As your brain processes each word, it updates its predictions, and, as each one pans out, the prediction and matching input cancel each other out. Because the grammatical sentence matches the brain's predicted structure, there is no significant prediction error, and processing is smooth.

But had I in fact come up with the less likely~-- in fact wildly unlikely~-- utterance \textit{the cat the milk drank}, there would be some updating to do. This mismatch between expectation and input would generate a prediction error. After all, you've estimated the probability of such and outcome to be vanishingly tiny. Now you have to work out what structure and meaning I had intended and decide whether something truly strange is going on with your cat, whether I've lost my grammatical faculties, whether you've misheard, or any number of other possibilities. This all plays out during a brief hesitation while you stare incredulously at me, and you conclude that what I meant to utter was simply the banal observation about the cat having drunk its milk. At the same time, you will experience a strong feeling of ungrammaticality, as what I said violates the standard structure for conveying such an idea.

This reflects the fact that the ungrammatical sentence, while interpretable, is a much poorer fit to your brain's linguistic model. The larger prediction errors and the lower final probability of the chosen structure reflect the brain's struggle to make sense of the ungrammatical input, and this is what we experience as the feeling of ungrammaticality.

\bigskip

This Bayesian perspective helps to make sense of the graded and context-sensitive nature of grammaticality judgments. A construction that is highly probable in one context may be highly improbable in another, leading to different levels of prediction error and different degrees of perceived ungrammaticality.

It also helps to explain how our sense of grammaticality can change over time, as the probabilities in the Bayesian equation are updated based on new linguistic experiences. A construction that initially feels ungrammatical may gradually become more acceptable as we encounter it more often, just as the Tristan chord went from sounding wrong to sounding fine as audiences became more familiar with it.

Of course, the brain isn't literally calculating probabilities and arriving at a precise likelihood. These Bayesian computations are implicit in the neural machinery of prediction and error correction. But the mathematical framework provides a useful way to think about the cognitive processes underlying our grammaticality judgments.

\subsection{Grammar sensitivity in animals}

The sensation of grammaticality we experience is often taken as evidence for the human-specific nature of language. But could this feeling have emerged from cognitive abilities that predate human language? Recent research with non-human primates suggests that the mental machinery underlying our sensitivity to grammatical dependencies may have evolved long before language itself.

In a groundbreaking study, \citet{Ravignani2013} demonstrated that squirrel monkeys (\textit{Saimiri sciureus})~-- New World primates whose lineage diverged from ours at least 36 million years ago~-- can detect abstract, non-adjacent dependencies in auditory sequences. The monkeys were habituated to tone sequences following an AB$^n$A pattern, where the first and last elements matched while being separated by a variable number of intervening elements. Not only did the monkeys distinguish between sequences containing dependencies and those lacking them, but they also generalized to previously unheard pitch classes and novel dependency distances.

This is remarkable precisely because squirrel monkeys lack vocal learning abilities and possess a relatively fixed vocal repertoire. These monkeys will never speak a human language, yet they possess the cognitive toolkit to recognize abstract patterns that form the foundation of many linguistic structures. As the authors note, ``Although no squirrel monkey will probably ever speak a human language, these monkeys possess the cognitive potential to recognize the rule generating plurals of Turkish nouns, or many other linguistic phenomena'' \citep[4]{Ravignani2013}.

What makes this finding so significant for our understanding of (un)gram\-maticality feelings is that it suggests a profound disconnect between the evolution of language and the evolution of dependency sensitivity. Rather than evolving specifically for linguistic purposes, our feeling of ungrammaticality when encountering a malformed sentence may be an \textsc{exaptation}~-- the repurposing of a cognitive ability that evolved for entirely different reasons.

Like the feeling of disgust that evolved to help us avoid pathogens but now extends to moral violations, our sense of grammaticality likely built upon pre-existing pattern-recognition systems that long predated human language. Our hominid ancestors may have used these skills for parsing complex environmental patterns, recognizing social structures, or even processing the temporal sequences in conspecific vocalizations, long before the emergence of syntactic language.

This perspective helps explain why our feelings of (un)grammaticality can be so context-dependent and variable. The neural circuitry underlying these feelings wasn't optimized specifically for language but was co-opted from more general-purpose cognitive mechanisms that had evolved for other adaptive functions. The mechanism~-- detecting violations of expected patterns~-- served our ancestors well in many domains of life, and language simply hijacked this pre-existing system.

The squirrel monkey study also emphasizes how these mechanisms are inherently predictive. The monkeys weren't simply memorizing specific sequences; they were extracting abstract patterns and using them to make predictions about novel stimuli. When these predictions were violated, the monkeys showed increased attention~-- the animal equivalent, perhaps, of our feeling of ungrammaticality. This aligns with our discussion of prediction error as a core component of grammaticality intuitions.

If a squirrel monkey can detect a violation in a tone sequence following the same abstract pattern as Hungarian vowel harmony, then our own feelings about well-formedness in language likely arose from deep cognitive foundations shared across primates. Our human language faculty didn't create sensitivity to dependencies; it inherited and specialized mechanisms that were already millions of years old.

This evolutionary perspective challenges the notion that grammar is inherently linguistic. Rather, grammar~-- and our feelings about it~-- may be better understood as the application of domain-general pattern-recognition abilities to the specific realm of language. What we experience as a feeling of ungrammaticality may be nothing more than the same surprise signal our ancestors felt when encountering any unexpected pattern in their environment, now directed at linguistic stimuli.

\subsection{The social role of ungrammaticality feelings}

If our sense of ungrammaticality arose before human language, what purpose might it have served? Evidence from non-human primates and other social mammals offers compelling clues. The studies on squirrel monkeys by \citet{Ravignani2013} and along with others on rock hyraxes by \citet{Kershenbaum2012} reveal that syntactic processing capabilities may have evolved for social purposes rather than for language per se.

Rock hyraxes, look like scruffy rodents but are actually more closely related to elephants and manatees~-- a biological surprise package about the size of a robust rabbit. These charming, vocal mammals live in colonies among rocky outcrops across Africa and the Middle East, where they spend their days sunbathing, foraging, and engaging in complex social interactions punctuated by elaborate songs. That such an unassuming creature produces vocalizations with sophisticated syntactic structure challenges our assumptions about which animals might possess complex communication systems.

In squirrel monkeys, the ability to detect abstract dependencies between non-adjacent elements suggests a cognitive capacity for recognizing patterns in vocal sequences. Similarly, hyraxes demonstrate complex syntactic structures with geographical dialects that correlate with social proximity. Both findings point to a potential social function of syntax detection: group recognition and cohesion.

What we experience as a feeling of ungrammaticality might have originated as a social signal~-- an alert to potential out-group communication. For our primate ancestors, detecting ``incorrect'' vocal patterns could help identify unfamiliar individuals or groups, triggering appropriate caution or interest. This mechanism would be particularly valuable in environments where visual identification was limited, such as dense forests or during low-light conditions.

The correlation between vocal similarity and geographical distance in rock hyraxes suggests that syntactic patterns may serve as markers of social familiarity. As \citet{Kershenbaum2012} observed, song syntax similarity degrades with distance, but only within a limited range corresponding to typical dispersal patterns. This pattern is consistent with a social signaling system that evolved to distinguish local group members from strangers.

Moreover, the mother-daughter vocal interactions in gibbons described by \citet{Koda2013} reveal how syntactic patterns might be socially transmitted. Their finding that ``co-singing interactions between mothers and daughters may enhance vocal development'' suggests that social learning plays a crucial role in acquiring group-specific syntactic patterns.

Our human sense of ungrammaticality, then, may be an evolutionary repurposing of a much older social recognition system. What was once a mechanism for distinguishing in-group from out-group vocalizations has become a sophisticated tool for parsing linguistic structure. But the visceral nature of our reactions to ungrammatical sentences~-- that feeling of wrongness~-- may reflect this deep evolutionary history.

This evolutionary heritage may help explain the surprisingly emotional responses many people have to ``grammatical errors'' in modern contexts. Language policing~-- the practice of publicly correcting others' speech or writing~-- rarely stems from genuine communication needs. Rather, it often functions as a social gatekeeping mechanism. The grammar pedant who bristles at a split infinitive or corrects someone's \textit{whom} usage may be unconsciously engaging an ancient social marking system, one that evolved to identify who belongs and who doesn't.

The strong negative reactions to dialectal differences, casual speech patterns, or linguistic innovations reveal how deeply intertwined our sense of grammaticality is with social identity. When we feel that visceral discomfort at hearing someone ``break the rules" of grammar, we may be experiencing the echo of an evolutionary adaptation that once helped our ancestors distinguish friend from stranger. In modern contexts, however, these instinctive reactions often serve to reinforce social hierarchies and exclude those who speak differently~-- suggesting that our feelings about grammaticality remain as much about social boundaries as they are about language itself.

\section{Grammar and the Sacred}

The Tower of Babel story captured something true about language. God didn't just divide people by making them speak differently~-- he divided them by making their grammars incompatible. Faced with patterns that felt profoundly wrong, they could no longer work together. Each group took their own grammatical patterns as natural and correct. The story tells us that grammatical differences are not mere technical distinctions but markers of fundamental social boundaries.

When sociologists talk about `the sacred', they mean patterns that bind communities together and mark them off from others. Sacred patterns feel natural and right to insiders but strange or wrong to outsiders. They persist through enormous social investment but feel effortless to practitioners. They change, but slowly and from within.

Grammar works this way. German puts many verbs at the end of clauses. English doesn't. Each pattern feels absolutely natural to its speakers. An English speaker encountering \textit{*I a walk take} recoils as spoken to by a creature from a swamp on the planet Dagobah. But German speakers do this naturally all day long. The pattern that feels profane to one community is sacred to another.

Southern American English shows this clearly in its double modals. Speakers say \textit{I might could help you} or \textit{He should oughta be here by now}. These constructions are completely natural within their speech communities. Children acquire them without effort. Standard English speakers find them deeply ungrammatical~-- not just incorrect but somehow wrong. Both reactions reveal grammar's sacred character.

The profane appears when sacred patterns are violated. Consider verb complements in English. We say \textit{decide to go} but \textit{enjoy going}. Mix these patterns up (\textit{*decide going}, \textit{*enjoy to go}) and English speakers react viscerally. The violations feel not just wrong but almost unclean. This mirrors how sacred traditions treat violations of fundamental patterns.

The sacred status of grammatical patterns explains why some changes succeed while others fail. When English needed new ways to talk about possession, \textit{his book} extended naturally to \textit{Mary's book}. The innovation grew from existing patterns. But conscious attempts to change grammar usually fail. Invented pronouns like \textit{ze} or \textit{xe} never caught on. They violated the sacred character of grammatical patterns by imposing new forms from outside.

The sacred appears most clearly in how communities maintain their distinct grammars despite close contact. African American English keeps its aspectual system because these patterns carry deep community meaning. Take the remote past marker \textit{been}, as in \textit{She been married}. Insiders process this automatically. Outsiders systematically misunderstand. The patterns persist not through conscious effort but through their role in marking sacred social boundaries.

These boundaries work in subtle ways. Indian English speakers say \textit{I am having two books} naturally. Many other English speakers find this wrong. Each group's patterns feel sacred through use. But `sacred' here doesn't mean unchangeable. Consider how English developed the going-to future (\textit{I'm going to leave}). The construction emerged gradually through speakers finding new ways to express meaning. Similar patterns appear across languages~-- French future tense from Latin \textit{habere}, Mandarin perfective \textit{-le} from `finish'. Each change preserved sacred patterns while adapting to new needs.

Grammar differs from many sacred traditions in working below conscious awareness. French authorities campaign for linguistic purity. But basic French grammar patterns persist through implicit learning. Children master complex syntax without explicit teaching. The patterns become sacred through use rather than doctrine.

Grammar's sacred character explains why grammatical intuitions feel both universal and local. When English speakers judge \textit{*Book the him gave she} as wrong, they're not consulting rules. Their reaction comes from deep patterns of community practice~-- patterns as automatic as any sacred tradition. But these patterns vary between communities because they emerge from shared histories of use.

The story of Babel got this right. Grammar divides us precisely because it unites us. The patterns that bind one community together mark boundaries with others. We experience violations of our grammatical patterns as profane because these patterns provide the sacred ground of shared meaning. Grammar is not sacred in any mystical sense. But it takes on sacred characteristics because it enables the fundamental human task of making meaning together.

\section{Barrett's Theory of Constructed Emotion: Implications for Grammaticality}

\subsection{Overview of the theory of constructed emotion}

Lisa Feldman Barrett's theory of constructed emotion offers a paradigm shift in how we understand emotions~-- and by extension, how we might understand other psychological phenomena, including our sense of grammaticality. Just as Barrett argues that emotions are not innate, universal categories but rather constructed experiences, we might consider whether our intuitions about grammaticality are similarly constructed rather than innate or universal.

Barrett's theory posits that emotions emerge from more basic psychological ingredients: interoception (the brain's representation of sensations from the body), concepts (knowledge of the world based on past experience), and affect (feelings of pleasantness/unpleasantness and arousal/calmness). The brain, she argues, is constantly predicting and checking its predictions against sensory input, using concepts to categorize sensations and create meaningful experiences~-- including emotions.

Consider how this might apply to grammaticality. When we encounter a sentence, our brain isn't passively receiving input. Instead, it's actively predicting what's coming next, based on our past linguistic experiences (our `concepts' of language). When we encounter something unexpected~-- say, a sentence like *\textit{The child seems sleeping}~-- our prediction is violated. This violation might generate an affective response (perhaps a feeling of unpleasantness or increased arousal), which we then categorize using our concept of `ungrammaticality'.

This view suggests that our sense of what's grammatical might be intimately tied to our broader cognitive and affective processes, constructed in the moment based on our linguistic experiences and current context.

\subsection{Key concepts: interoception, concepts, and affect}

To fully appreciate the implications of Barrett's theory for our understanding of grammaticality, we need to delve deeper into its key components: interoception, concepts, and affect.

\textbf{Interoception} refers to the brain's representation of the body's internal state. It's not just about conscious awareness of heartbeats or digestion, but a constant, usually unconscious process of monitoring and predicting the body's needs. In the context of language processing, we might consider how our physiological state influences our perception of language. Are we more likely to judge a sentence as ungrammatical when we're tired or stressed? Does the effort of processing a complex sentence structure create bodily sensations that we interpret as signs of ungrammaticality?

\textbf{Concepts}, in Barrett's framework, are not static, dictionary-like definitions, but flexible, context-dependent predictions about what sensations mean. They're built from past experiences and constantly updated. Applied to grammar, this suggests our concepts of `correct' language are shaped by every linguistic interaction we've had, from formal education to casual conversation. When we judge a sentence like \textsuperscript{?}{\textit{The data is conclusive}} (some people prefer \textit{are}), we're not consulting an internal rulebook, but rapidly constructing a prediction based on our accumulated experiences with similar phrases.

\textbf{Affect} is the basic sense of feeling good or bad, calm or agitated. It's the currency of allostasis~-- our brain's process of anticipating and meeting the body's needs. In grammaticality judgments, affect might manifest as the subtle feelings of rightness or wrongness we experience when reading or hearing language. The discomfort we feel at a `garden path' sentence like \textit{The horse raced past the barn fell} might be a form of affective response, signaling the need to reallocate cognitive resources.

Together, these components form a dynamic system. Our brain predicts what sensations (including linguistic input) mean for our body budget, categorizes these sensations using concepts built from past experience, and generates affect to guide action~-- all in service of maintaining allostasis. This process, Barrett argues, constructs our emotional experiences. And perhaps, it also constructs our experience of grammaticality.

\subsection{Parallels between emotion construction and grammaticality judgments}

The parallels between Barrett's theory of constructed emotion and our experience of grammaticality are striking and potentially illuminating.

First, consider the variability of both emotions and grammaticality judgments. Just as the experience of `fear' can vary widely across individuals and contexts, so too can judgments of grammaticality. What feels ungrammatical to a copy editor might feel perfectly natural to a casual speaker. This variability is difficult to explain if we assume innate, universal categories of either emotions or grammatical rules, but it aligns well with a constructionist view.

Second, both emotions and grammaticality judgments involve a blending of `bottom-up' sensory input and `top-down' conceptual knowledge. In emotion, we interpret bodily sensations through the lens of our emotional concepts. Similarly, in language processing, we interpret strings of words through the lens of our grammatical knowledge. A sentence like \textit{The old man the phones} feels ungrammatical at first because our initial prediction (based on frequent patterns) is violated, but can feel grammatical once we apply different conceptual knowledge.

Third, both processes appear to be deeply contextual. The same physiological state might be conceptualized as `fear' or `excitement' depending on the context; similarly, the same linguistic construction might be judged grammatical or ungrammatical depending on the context. Consider how our judgments of sentences like \textit{Who do you wanna go to the movies with?} can shift depending on whether we're in a formal or casual setting.

Fourth, both emotional and grammatical concepts seem to be learned and culturally influenced. Just as emotional concepts can vary across cultures, so too can concepts of grammaticality. What's considered ungrammatical in Standard American English might be perfectly grammatical in another dialect, for instance.

Finally, both experiences of emotions and grammaticality judgments seem to involve an element of affect~-- a feeling of goodness or badness, rightness or wrongness. The `gut feeling' we get about whether a sentence is grammatical might be analogous to the core affect that Barrett posits as a basic ingredient of emotion.

These parallels suggest that our experience of grammaticality might be constructed in much the same way as our experience of emotion: through a dynamic interplay of sensory input, conceptual knowledge, and affective feeling, all in service of making meaning and guiding action in our linguistic world.

\subsection{Reinterpreting grammatical intuitions through a constructionist lens}

If we take seriously the idea that grammaticality judgments are constructed in a manner similar to emotions, we're led to some intriguing reinterpretations of familiar linguistic phenomena.

Consider garden path sentences like \textit{The horse raced past the barn fell}. The traditional view might explain our initial difficulty with such sentences in terms of parsing mechanisms or structural ambiguity. A constructionist view, however, might see this as a case of prediction error. Our brain, drawing on past linguistic experience, makes a strong prediction about the structure of the sentence after encountering \textit{The horse raced past the barn}. When \textit{fell} appears, violating this prediction, we experience a momentary increase in arousal and perhaps a feeling of unpleasantness~-- core affect that we then conceptualize as `ungrammaticality'.

Or consider our varying intuitions about sentences like \textit{The data is conclusive} versus \textit{The data are conclusive}. Rather than appealing to competing grammatical rules, we might view this as a case of competing predictions, each associated with different conceptual knowledge (formal academic writing vs. everyday usage) and perhaps different affective states (the smugness of formal correctness vs. the comfort of colloquial familiarity).

The phenomenon of \textsc{syntactic satiation}~-- where repeatedly reading an ungrammatical sentence makes it begin to feel grammatical~-- takes on new significance in this light. Rather than a fatigue of grammatical judgment faculties, we might see this as a rapid updating of our predictive model. Each exposure adjusts our concept of what's linguistically possible, gradually shifting our affective response from `wrong' to `right'.

Even the intuitions of trained linguists, often treated as more reliable than those of naive speakers, might be reinterpreted. Perhaps what distinguishes the linguist's judgments is not access to a more accurate internal grammar, but a richer set of linguistic concepts allowing for more nuanced prediction and categorization of language input.

This constructionist view also offers a new perspective on linguistic prescriptivism. The strong feelings that prescriptivists have about `correct' language use might be understood not as defense of innate grammatical rules, but as affective responses constructed from deeply ingrained concepts about language, social status, and cognitive sophistication.

By reframing grammatical intuitions as constructed experiences rather than readouts of an internal rule system, we open up new ways of understanding the flexibility, context-sensitivity, and sometimes puzzling inconsistency of our linguistic judgments.

\subsection{Reinterpreting grammatical intuitions through a constructionist lens}

If we take seriously the idea that grammaticality judgments are constructed in a manner similar to emotions, we're led to some intriguing reinterpretations of familiar linguistic phenomena.

Consider garden path sentences like \textit{The horse raced past the barn fell}. The traditional view might explain our initial difficulty with such sentences in terms of parsing mechanisms or structural ambiguity. A constructionist view, however, might see this as a case of prediction error. Our brain, drawing on past linguistic experience, makes a strong prediction about the structure of the sentence after encountering \textit{The horse raced past the barn}. When \textit{fell} appears, violating this prediction, we experience a momentary increase in arousal and perhaps a feeling of unpleasantness~-- core affect that we then conceptualize as `ungrammaticality'.

Or consider our varying intuitions about sentences like \textit{The data is conclusive} versus \textit{The data are conclusive}. Rather than appealing to competing grammatical rules, we might view this as a case of competing predictions, each associated with different conceptual knowledge (formal academic writing vs. everyday usage) and perhaps different affective states (the smugness of formal correctness vs. the comfort of colloquial familiarity).

The phenomenon of \textit{syntactic satiation}~-- where repeatedly reading an ungrammatical sentence makes it begin to feel grammatical~-- takes on new significance in this light. Rather than a fatigue of grammatical judgment faculties, we might see this as a rapid updating of our predictive model. Each exposure adjusts our concept of what's linguistically possible, gradually shifting our affective response from `wrong' to `right'.

Even the intuitions of trained linguists, often treated as more reliable than those of naive speakers, might be reinterpreted. Perhaps what distinguishes the linguist's judgments is not access to a more accurate internal grammar, but a richer set of linguistic concepts allowing for more nuanced prediction and categorization of language input.

This constructionist view also offers a new perspective on linguistic prescriptivism. The strong feelings that prescriptivists have about `correct' language use might be understood not as defense of innate grammatical rules, but as affective responses constructed from deeply ingrained concepts about language, social status, and cognitive sophistication.

By reframing grammatical intuitions as constructed experiences rather than readouts of an internal rule system, we open up new ways of understanding the flexibility, context-sensitivity, and sometimes puzzling inconsistency of our linguistic judgments.

\section{A neuroscientific perspective}

Ever catch yourself finishing someone else's sentence? That's your brain doing what it does best~-- making predictions. Our brains aren't just passive receivers of language; they're constantly trying to guess what's coming next. When those guesses are wrong, it's like a little alarm goes off in our heads. This predictive processing is key to understanding how our brains handle grammar.

Here's a fun example. Read this sentence:

\ea{\textit{How many animals of each kind did Moses take on the Ark?}}\label{ex:moses}
\z

Did you spot the error? It was Noah, not Moses. But don't feel bad if you missed it~-- most people do. That's because your brain, making predictions based on the words ``animals'', ``Ark'', and a biblical figure, fills in the blanks so efficiently that it often overrides your actual knowledge. It's like your brain is saying, ``Close enough!''

This predictive power explains why we can breeze past some pretty glaring grammatical errors. Take this gem from \textit{Fifty Shades of Grey}:

\ea[*]{\textit{Pulling off his boxer briefs, his erection springs free. Holy cow!}}\label{ex:erection}
\z

Grammatically, this is a howler. \textit{Pulling off his boxer briefs} should describe what the subject of the sentence is doing, but unless this is a very different kind of novel, \textit{his erection} isn't doing any undressing. Yet most readers don't even blink. The broader context creates such strong expectations that our brains just fill in the blanks and move on.

Our brains don't just passively receive language input~-- they're constantly making predictions about what's coming next. When those predictions are violated, it generates a neural surprise signal. This predictive processing framework offers a powerful lens for understanding how our brains handle grammar.

Recent research (doi.org/10.1038/s41586-024-07973-1) has given us an unprecedented look at how this process works at the level of individual neurons. In a groundbreaking study, scientists recorded the activity of single neurons in the brains of epilepsy patients who had electrodes implanted for clinical reasons. They found that neurons in the hippocampus and entorhinal cortex~-- areas crucial for memory and navigation~-- gradually changed their firing patterns to represent the structure of a sequence of images the patients were viewing.

This neuronal representation formed rapidly and without any explicit instructions. The patients weren't told to look for patterns; their brains just naturally extracted the underlying structure. This is remarkably similar to how we learn grammar~-- not through memorizing rules, but by implicitly picking up on patterns in the language we hear.

Even more intriguingly, when the researchers looked at the combined activity of many neurons, they could recover the entire structure of the image sequence. It's as if the brain had built a mental map of the sequence, much like the ``cognitive maps'' these same brain areas create of physical spaces. This suggests that our brains might represent the structure of language~-- its grammar~-- in a similar way to how we represent the layout of a city or the order of events in a story.

This neuronal representation wasn't just a passive record of what had happened. It was predictive, encoding information about what was likely to come next in the sequence. This aligns perfectly with the predictive processing framework we discussed earlier. Your brain isn't just recognizing grammatical sentences; it's actively predicting how those sentences will unfold.

The researchers also found evidence of ``replay''~-- the neurons would spontaneously reactivate in the same patterns they'd followed during the initial sequence, but speeded up. This replay happened during breaks, suggesting it might be a way for the brain to consolidate what it had learned. It's not hard to imagine similar processes happening as we internalize the patterns of our native language.

Let's bring this back to grammar. When you encounter a sentence like \textit{The old man the phones}, your brain initially predicts that \textit{man} is a noun, based on the familiar pattern of \textit{The old man...} But as soon as you hit \textit{the boats}, that prediction is violated. Your neurons are likely firing in patterns that reflect this violation, similar to the ``surprise'' signals seen in the image sequence study. Your brain then has to rapidly revise its interpretation, now understanding \textit{man} as a verb.

Neuroscientists have actually measured these prediction processes in the brain. When we encounter unexpected words or structures, our brains generate distinct patterns of electrical activity. There are two main ones to know about:

\begin{itemize}
    \item The N400: Think of this as your brain's ``semantic speed bump''. It happens about 400 milliseconds after you encounter a word that doesn't fit semantically. If I say \textit{I like my coffee with cream and socks}, your brain hits that speed bump at ``socks''.
    \item The P600: This is more like a ``syntactic pothole''. It occurs about 600 milliseconds after you hit a grammatical snag. You'd get a big P600 if you read \textit{The man bit dog the}.
\end{itemize}

These brain blips get bigger the more surprising the word or structure is. A mild surprise is like a gentle speed bump, while a major violation is like hitting a pothole at 60 mph.

But here's where it gets really interesting: context can totally change how our brains respond. That Moses sentence earlier? In casual conversation, your brain might just sail right past it without triggering an N400 response. But in a psychology experiment, where you're primed to expect tricky questions and are paying close attention, you're more likely to notice the error. In that case, your brain would probably generate an N400 response when it encounters \textit{Moses}.

This flexibility ties into a broader idea in neuroscience: that our experiences, including our sense of what's grammatical, might be actively constructed by our brains. Lisa Feldman Barrett proposed this for emotions~-- that they're not fixed categories, but experiences our brains construct on the fly from more basic ingredients.

Kent Berridge extended this idea to our sense of good and bad~-- what neuroscientists call ``valence''. He challenged the idea that specific brain regions are permanently dedicated to positive or negative feelings. Instead, he suggested that neural systems can have multiple ``affective modes'', capable of producing both positive and negative feelings depending on the context.

We can apply this same thinking to grammaticality. When we encounter a sentence, we're not just mechanically applying rules. We're rapidly constructing an experience based on:

\begin{itemize}
    \item How we're feeling physically (Are we tired? Stressed?)
    \item What we know about language
    \item Our gut feeling of ``rightness'' or ``wrongness''
\end{itemize}

This explains why our sense of what's grammatical can be so flexible. Just like how the same physical sensation might feel like excitement or anxiety depending on the context, the same sentence structure might feel grammatical or ungrammatical depending on the situation.

Take a sentence like \textit{That's the person who I saw}. In formal writing, you might get a twinge of wrongness~-- shouldn't it be ``whom''? But in casual conversation, it'll likely sail right by. Your brain, in different contexts, is in different ``grammatical modes''.

This view also sheds light on the phenomenon of syntactic satiation~-- where if you read an ``ungrammatical'' sentence enough times, it starts to feel okay. You're not wearing out your grammar module; you're updating your brain's predictive model. What was once surprising becomes expected.

It even explains why linguists often have different grammatical intuitions than the general public. It's not that linguists have access to a more accurate internal grammar. They've just developed a richer set of linguistic concepts, allowing for more nuanced predictions and categorizations.

This neuroscientific perspective invites us to see grammaticality not as a binary property of sentences, but as a dynamic, constructed experience. It arises from the interplay of prediction, sensation, and conceptualization. This view embraces the messiness and flexibility of real-world language use, while still grounding our intuitions in the fundamental workings of our predictive brains.

Berridge's work provides some concrete examples of this flexibility. He found that stimulating neurons in the central amygdala~-- an area often associated with fear~-- could actually enhance positive motivation for rewards like cocaine or sugar. Similarly, in the nucleus accumbens, the same neural populations could drive both positive (eating) and negative (defensive) behaviors depending on the environment.

If we apply this to grammaticality, it suggests our neural circuits for language processing might be equally flexible. The same brain areas that spark a sense of wrongness when you say \textit{Me and him went to the store} might, in a different context, give you a pleasant frisson when a poet breaks the rules for artistic effect.

This doesn't mean grammatical rules don't exist or don't matter. But it does suggest that our experience of those rules~-- our feeling of grammaticality~-- is far more dynamic and context-dependent than we might have thought. It's not about matching input to a stored set of rules, but about constructing a linguistic experience that best fits our current model of the world.

In this light, the feeling of ungrammaticality becomes less about detecting a rule violation and more about experiencing a prediction error. And just as our brains can rapidly update their predictions in other domains, they can also update their linguistic predictions. This is how new constructions gradually become acceptable, and how the language evolves over time.

This neuroscientific view offers a richer, more nuanced understanding of grammaticality. It helps explain why our intuitions can sometimes feel fuzzy or contradictory, why context matters so much, and why there's often not a clear line between ``grammatical'' and ``ungrammatical''. It suggests that our sense of grammar is not a fixed, innate system, but a dynamic, constructed experience~-- one that's intimately tied to our broader cognitive and emotional processes.

\section{Form-meaning mismatch}

Open a dictionary to any random page to check the definition of a word, and you'll rarely find one~-- instead, you'll typically be presented with a list of senses. Take, for example, the word \textit{bank}. The \textit{Oxford English Dictionary} lists three distinct entries for \textit{bank} as a noun alone: `the edge of a river', `an institution that invests money', and `a set or array of similar items'. Yet this only scratches the surface of its semantic complexity. Even within these broad categories, numerous sub-senses exist. The `edge of a river' definition, for instance, is one of two major senses, alongside the earlier meaning of `a mound or slope'. Together, these two broad senses encompass a total of twelve sub-senses. 

The same multiplicity of meanings~-- or \textsc{polysemy}~-- can be found in grammar. In my model, I claim that, at its most basic, a grammatical construction is one with an accepted form–meaning pairing. But I also note that a grammatical construction may evolve additional accepted form–meaning pairings.

But both in words and grammar, we sometimes act as though everything should be monosemous~-- it should have only a single meaning. The other meanings, the scolds will tell us, are errors. Take the case of \textit{can}. Far too many children have been tormented by teachers who will not accept \textit{can I go to the washroom} but insist on \textit{may}. This stems from the teachers' wrongheaded belief that \textit{can} has the meaning of ability, while \textit{may} has the meaning of permission. Any dictionary will belie this fallacious position.

When it comes to grammar, the fallacy of monosemy can be even more entrenched because of the names we give to grammatical systems. Terms like \textsc{the past tense}, \textsc{singular}, \textit{definite}, etc. give us tunnel vision about the possible meanings of the forms in question.

Certainly. I'll build out the two examples at the end, expanding on the concepts of the past tense and definiteness to illustrate the form-meaning mismatch in grammar. Here's an improved version:
latexCopy\section{Form-meaning mismatch}
Open a dictionary to any random page to check the definition of a word, and you'll rarely find one~-- instead, you'll typically be presented with multiple meanings. Take, for example, the word \textit{bank}. The \textit{Oxford English Dictionary} lists three distinct entries for \textit{bank} as a noun alone: `the edge of a river', `an institution that invests money', and `a set or array of similar items'. Yet this only scratches the surface of its semantic complexity. Even within these broad categories, numerous sub-senses exist. The `edge of a river' definition, for instance, is one of two major senses, alongside the earlier meaning of `a mound or slope'. Together, these two broad senses encompass a total of twelve sub-senses. 

The same multiplicity of meanings~-- or \textsc{polysemy}~-- can be found in grammar. In my model, I claim that, at its most basic, a grammatical construction is one with an accepted form--meaning pairing. But I also note that a grammatical construction may evolve additional accepted form--meaning pairings. But this can be counter-intuitive.

Both in words and grammar, our intuition often suggests to us that everything should be monosemous~-- it should have only a single meaning. The other meanings, the scolds will tell us, are errors.

Take the case of \textit{can}. Far too many children have been tormented by teachers who will not accept \textit{can I go to the washroom} but insist on \textit{may}. This stems from the teachers' wrongheaded belief that \textit{can} has the meaning of ability, while \textit{may} has the meaning of permission. Any dictionary will belie this fallacious position. As will some casual reading or observation of conversations. But the intuitive pull of the fallacy of monosemy is strong.

When it comes to grammar, belief in monosemy can be even more entrenched, likely because of the names we give to grammatical systems. Terms like \textsc{the past tense}, \textsc{singular}, \textit{definite}, etc. give us tunnel vision about the possible meanings of the forms in question. Clearly, the past tense should be fore past time, singular nouns should denote individuals, and a definite NP should be specific and identifiable. But grammatical systems, like words, tend to be polysemous.

Consider the past tense. While its primary function is indeed to locate events in the past, it serves several other purposes in English:

\begin{enumerate}
    \item Hypothetical situations: \textit{If I won the lottery, I would buy a house.}
    \item Politeness: \textit{I wanted to ask you a favour.}
    \item Reported speech: \textit{She said she was busy.}
    \item Narrative style in storytelling: \textit{Once upon a time, there was a princess...}
\end{enumerate}

Each of these uses demonstrates that the past tense form carries meanings beyond simple past time reference, illustrating the mismatch between its name and its full range of functions.

Similarly, definiteness, typically associated with the definite article `the', exhibits a range of meanings:

\begin{enumerate}
    \item Unique reference: \textit{The sun is shining.}
    \item Shared knowledge: \textit{Did you feed the cat?}(assuming both speakers know which cat)
    \item Generic reference: \textit{The lion is a majestic animal.}
    \item Superlatives: \textit{This is the best restaurant in town.}
    \item Unknown characters in anecdotes: \textit{I saw this guy.}
\end{enumerate}

Here, we see that definiteness encompasses not just a single concept of `being definite', but a range of semantic and pragmatic functions related to specificity, shared knowledge, and even generality.

Innocent people can thus be lulled into persuading themselves that something is ungrammatical even when it doesn't feel ungrammatical. At least, not initially.

\section{One intuition or many}
Joshua Knobe (forthcoming from Ergo) argues that not only do different people have different intuitions but that individuals have conflicting intuitions.

The following is from chatGPT:

Modeling the brain as a system that calculates vectors within a high-dimensional vector space provides a powerful framework for understanding complex phenomena like conflicting emotions. Here's how conflicting emotions might manifest within this model:
1. Representation of Emotions as Vectors

    Dimensionality: Each emotion can be represented as a vector in a high-dimensional space, where each dimension corresponds to different emotional attributes or intensities. For instance, dimensions could capture aspects like valence (positive vs. negative), arousal (calm vs. excited), or specific emotions (joy, sadness, anger, etc.).

    Vector Components: The position and direction of an emotion vector are determined by its components along these dimensions. A specific emotional state is thus a point in this high-dimensional space, characterized by the combination of active emotions.

2. Conflicting Emotions as Divergent Vectors

    Opposing Directions: Conflicting emotions arise when multiple emotion vectors point in divergent or even opposing directions within the vector space. For example, feeling both happiness and sadness about a particular event would involve two vectors~-- one representing happiness and the other sadness~-- that are not aligned and may partially or wholly oppose each other.

    Vector Interference: When these conflicting vectors are combined (e.g., through vector addition), they can interfere with each other, leading to a resultant vector that doesn't fully capture either emotion. This interference can create a sense of emotional tension or ambiguity, as the brain simultaneously processes multiple, incompatible emotional signals.

3. Cognitive Processing of Conflicting Emotions

    Dynamic Interactions: In a high-dimensional space, the brain continuously processes and updates these emotional vectors. Conflicting emotions might lead to dynamic shifts as different vectors gain or lose influence based on context, memories, or external stimuli.

    Resolution Mechanisms: The brain may employ mechanisms to resolve these conflicts, such as prioritizing certain emotions over others or integrating aspects of each emotion to form a more nuanced emotional state. In vector terms, this could involve adjusting the magnitude or direction of specific vectors to achieve a more harmonious overall state.

4. Implications for Behavior and Decision-Making

    Complex Decision Landscapes: Conflicting emotions create complex landscapes in the vector space, influencing decision-making processes. The resultant emotional vector impacts motivations, preferences, and choices, often requiring the individual to navigate through the tension between opposing emotional forces.

    Neural Computation: Understanding emotions as vectors allows for the exploration of how neural computations balance these conflicting signals, potentially shedding light on phenomena like cognitive dissonance, emotional regulation, and the subjective experience of mixed emotions.

5. Practical Applications

    Psychological Modeling: This vector-based approach can aid in creating more accurate psychological models that account for the complexity of human emotions, including their simultaneous and conflicting nature.

    Artificial Intelligence: In AI and machine learning, such models can improve emotion recognition and generation systems, enabling more sophisticated and human-like emotional interactions.

Illustrative Example

Imagine someone receiving both praise and criticism for a project they worked hard on:

    Praise Vector: Points toward positive valence and high self-esteem.
    Criticism Vector: Points toward negative valence and self-doubt.

These two vectors coexist in the individual's emotional space, creating a resultant vector that reflects a mixture of pride and uncertainty. The individual experiences the tension between feeling accomplished and questioning their performance, influencing their subsequent actions and emotional state.
Conclusion

Modeling the brain's emotional processing as vector calculations in a high-dimensional space offers a nuanced understanding of how conflicting emotions coexist and interact. This framework captures the simultaneous activation of multiple, directionally diverse emotional states, illustrating the complexity of human emotional experiences and providing a foundation for further exploration in both neuroscience and artificial intelligence.